\exercisesection{正则局部环}

\begin{enumerate}
\def\labelenumi{\arabic{enumi})}
\item
  设 A 为一个环，S 为 A 的一个乘法子集。证明有 \(\mathrm{dh}(\mathrm{S}^{-1}\mathrm{A}) \leqslant \mathrm{dh}(\mathrm{A})\)（A, X, p.138, 定义 2）。（取 \(S^{-1}A\)-模的投射 A-模解并将其局部化。）
\item
  设 A 为诺特环，B 为忠实平坦的 A-代数。则有 \(\mathrm{dh(A)}\leqslant \mathrm{dh(B)}\)。
\end{enumerate}

3) 设 A 为诺特局部环。我们拟证明，A 为正则环当且仅当 dh(A) 有限。

\begin{enumerate}
\def\labelenumi{\alph{enumi})}
\tightlist
\item
  假设 dh(A) 有限。为证明 A 为正则环，使用 A, X, p.208, 习题 13。
\item
  假设 A 为正则环。设 x 是生成 m\_A 的完全相交族。则与 x 相关的 Koszul 复形是由自由 A-模构成的 A/m\_A 的有限解。然后应用 A, X, p.203, 习题 16。
\end{enumerate}

\begin{enumerate}
\def\labelenumi{\arabic{enumi})}
\setcounter{enumi}{3}
\item
  设 A 为正则诺特局部环。则有 \(\dim(A) = \mathrm{dh}(A)\)（为证明 \(\mathrm{dh}(A) \leqslant \dim(A)\)，使用与参数系相关的 Koszul 复形。然后确定群 \(\operatorname{Ext}_{A}^{i}(\kappa_{A}, \kappa_{A})\)。）
\item
  设 A 为正则诺特局部环，p 为 A 的一个素理想。则 \(A_{p}\) 为正则环（使用习题 1 和习题 3）。
\item
  若对于 A 的每个素理想 p，局部环 \(A_{p}\) 都是正则环，则称诺特环 A 为正则环。
\end{enumerate}

\begin{enumerate}
\def\labelenumi{\alph{enumi})}
\item
  设 A 为诺特环。则 A 为正则环当且仅当对于 A 的每个极大理想 m，\(A_{m}\) 都是正则环（使用习题 5）。
\item
  设 A 为诺特环。证明 A 为有限维正则环当且仅当 dh(A) 有限（使用习题 4）。
\item
  证明 p.83 习题 13 中描述的环是无限维正则环。
\item
  设 A 为正则诺特环，S 为一个乘法子集。证明 \(\mathbf{S}^{-1}\mathbf{A}\) 为正则环。
\item
  设 A 为正则诺特环。证明 A{[}X{]} 为正则环。（通过局部化归约到 dh(A) 有限的情形。使用 A, X, p. 143, 定理 1 及 b)。）
\end{enumerate}

\(^{1}\) 关于此习题的更多细节，可参阅 R. Stanley，《分次代数的 Hilbert 函数》，Advances in Math., 28 (1978), 第 57–83 页。

\begin{enumerate}
\def\labelenumi{\arabic{enumi})}
\setcounter{enumi}{6}
\tightlist
\item
  设 I 为含有可数共尾子集的定向集，A 为一个环（不必交换）。我们拟证明，以 I 为指标的投射（左）A-模的每个归纳极限的投射维数都 \(\leqslant 1\)（A, X, p.~134）。
\end{enumerate}

\begin{enumerate}
\def\labelenumi{\alph{enumi})}
\tightlist
\item
  设 \(((\mathbf{Q}_i)_{i\in \mathbf{I}},(\varphi_{i,j})_{i < j})\) 为投射 A-模的归纳系统，且 \(Q = \varinjlim_{I} Q_i\)。为证明 \(Q\) 的投射维数 \(\leqslant 1\)，归约到 \(I = \mathbf{N}\) 的情形。此时有正合序列
\end{enumerate}

\[
0 \longrightarrow \bigoplus_ {n \in \mathbf {N}} Q _ {n} \stackrel {u} {\longrightarrow} \bigoplus_ {n \in \mathbf {N}} Q _ {n} \longrightarrow Q \longrightarrow 0
\]

其中对于每个 \(n \in \mathbf{N}\) 及每个 \(x \in \mathbf{Q}_n\)，有

\[
u (x) = \varphi_ {n, n + 1} (x) - x.
\]

\begin{enumerate}
\def\labelenumi{\alph{enumi})}
\setcounter{enumi}{1}
\tightlist
\item
  设 \(((\mathbf{P}_i)_{i\in I}, (\varphi_{i,j})_{i < j})\) 为投射 A-模的归纳系统，使得对于每个对 \((i,j)\)（其中 \(i < j\)），\(\varphi_{i,j}\) 都是单射，且模 \(\operatorname{Coker}(\varphi_{i,j})\) 是投射的。证明模 \(\mathbf{P} = \varinjlim_{I} \mathbf{P}_i\) 是投射的。（设 \(0 \to \mathbf{M}' \to \mathbf{M} \to \mathbf{M}'' \to 0\) 是 A-模的正合序列。则
\end{enumerate}

\[
0 \to \operatorname{Hom} _ {\mathbf {A}} (\mathrm{P} _ {i}, \mathrm{M} ^ {\prime}) \to \operatorname{Hom} _ {\mathbf {A}} (\mathrm{P} _ {i}, \mathrm{M}) \to \operatorname{Hom} _ {\mathbf {A}} (\mathrm{P} _ {i}, \mathrm{M} ^ {\prime \prime}) \to 0
\]

这是交换群的射影系统的正合序列，且射影系统 \(i \mapsto \mathrm{Hom}_{\mathbf{A}}(\mathbf{P}_{i}, \mathbf{M}')\) 对离散拓扑具有 Mittag-Leffler 性质（TG, II, p.18）。由此可推出射影极限序列是正合的。）

\begin{enumerate}
\def\labelenumi{\alph{enumi})}
\setcounter{enumi}{2}
\tightlist
\item
  设 \((\mathbf{Q}_i)_{i\in I}\) 为投射 A-模的归纳系统。对于每个 \(i\in \mathbf{I}\)，令 \(\mathbf{R}_i = \bigoplus_{j\leqslant i}\mathbf{Q}_j\)，并且对于 \(i < i'\)，令 \(\psi_{i,i'}: \mathbf{R}_i \to \mathbf{R}_{i'}\) 为典范单射。定义归纳系统的正合序列 \(0 \to (\mathbf{P}_i) \to (\mathbf{R}_i) \to (\mathbf{Q}_i) \to 0\)，使得 \((\mathbf{P}_i)\) 满足 \(b\) 的假设。由此推出 \(a\) 的另一种证明。）
\end{enumerate}

\begin{enumerate}
\def\labelenumi{\arabic{enumi})}
\setcounter{enumi}{7}
\tightlist
\item
  设 A 为一个环，n 为满足 n ≥ 1 的整数。证明下列条件等价：
\end{enumerate}

\begin{enumerate}
\def\labelenumi{(\roman{enumi})}
\item
  有 dh(A) ≤ n。
\item
  对每个循环 A-模 M，都有 \(\mathrm{dp}_{\mathbf{A}}(\mathbf{M}) \leqslant n\)。
\item
  对 A 的每个理想 I，都有 \(\mathrm{dp}_{\mathbf{A}}(\mathbf{I}) \leqslant n - 1\)。
\end{enumerate}

(使用 A, X, p.138, 命题 4 证明等价性 (i) \(\Leftrightarrow\) (ii)，并使用 A, X, p.93, 命题 11 证明 (i) \(\Leftrightarrow\) (iii)。)

\begin{enumerate}
\def\labelenumi{\arabic{enumi})}
\setcounter{enumi}{8}
\tightlist
\item
  设 \(\Gamma\) 为全序交换群，\(\Gamma^{+}\) 为 \(\Gamma\) 的正元素集。我们以 (GOD) 表示下列性质：
\end{enumerate}

(GOD)：每个主子集 M ⊂ Γ⁺（VI, § 3, no. 5, 定义 2）都有一个关于逆序共尾的可数子集。

\begin{enumerate}
\def\labelenumi{\alph{enumi})}
\item
  证明若 \(\Gamma\) 的高度为 1，则满足性质 (GOD)。
\item
  对每个整数 n，证明存在一个具有性质 (GOD) 且高度为 n 的全序群 \(\Gamma\)（取赋予字典序的 \(\Gamma = \mathbf{Z}^{n}\)）。
\end{enumerate}

10) 设 A 为赋值环，其值的有序群 \(\Gamma\) 具有性质 (GOD)。a) 证明 A 的同调维数至多为 2。（注意 A 的每个理想都是主理想的可数归纳极限，并使用习题 7 和 8 的结果。）

\begin{enumerate}
\def\labelenumi{\alph{enumi})}
\setcounter{enumi}{1}
\item
  每个 Krull 维数为 1 的赋值环的同调维数都 \(\leqslant 2\)。它的同调维数为 2 当且仅当它不是诺特环。（使用习题 9, a) 以及 A, X, p.208, 习题 12。）
\item
  对每个整数 \(n \geqslant 1\)，存在 Krull 维数等于 n 且同调维数等于 2 的环。（使用习题 9, b) 以及 A, X, p.208, 习题 12。）
\end{enumerate}

\begin{enumerate}
\def\labelenumi{\arabic{enumi})}
\setcounter{enumi}{10}
\item
  每个正则诺特局部环都是唯一分解环。（使用 VII, § 4, n° 7, 推论 3。）
\item
  设 A 为一个环，a 为 A 的有限生成理想。
\end{enumerate}

\begin{enumerate}
\def\labelenumi{\alph{enumi})}
\item
  证明若 A 存在理想 b 使得 a = a·b，则存在 \(b \in b\) 使得 \((1 - b)a = (0)\)。
\item
  证明若 \(a^{2} = a\)，则存在元素 \(a \in a\) 使得 \(a^{2} = a\) 且 a = Aa。
\item
  假设 a 是极大理想，且 A-模 \(a/a^{2}\) 可由 n 个元素生成。证明理想 a 可由 \(n + 1\) 个元素生成。（取 a 中的元素 \(x_{1}, \ldots, x_{n}\)，它们在 \(a/a^{2}\) 中的像生成后者。记 x 为由 \(\{x_{1}, \ldots, x_{n}\}\) 生成的 A 的理想；注意在 A/x 中有 \((\mathfrak{a}/\mathfrak{x})^{2} = (\mathfrak{a}/\mathfrak{x})\)。）
\end{enumerate}

\begin{enumerate}
\def\labelenumi{\arabic{enumi})}
\setcounter{enumi}{12}
\tightlist
\item
  设 \(p\) 为素数，\(f\) 为满足 \(> 0\) 的整数。置 \(L = \mathbf{Q}(\zeta)\)，其中 \(\zeta\) 是 1 的一个本原 \(p^f\) 次根，并令 \(B = \mathbf{Z}[\zeta]\) 为 L 中由 \(\zeta\) 生成的子环。证明 B 是 Z 在 L 中的整闭包。（通过判别式计算可证明，\(B\left[\frac{1}{p}\right]\) 是 \(\mathbf{Z}\left[\frac{1}{p}\right]\) 的整闭包（V, § 1, no. 6, 引理 3）；记 \(\overline{\mathbf{B}}\) 为 \(\mathbf{B}\) 的整闭包，
\end{enumerate}

由此可推出 \(\overline{\mathbf{B}}\subset\bigcap_{q}\mathbf{Z}_{q}[\zeta]\)，其中 \(q\) 遍历素数，从而 \(\overline{\mathbf{B}}\subset\mathbf{Z}[\zeta].\)

\begin{enumerate}
\def\labelenumi{\arabic{enumi})}
\setcounter{enumi}{13}
\tightlist
\item
  设 \(\rho: A \to B\) 为诺特环的同态，使 \(B\) 成为忠实平坦的 \(A\)-模。
\end{enumerate}

\begin{enumerate}
\def\labelenumi{\alph{enumi})}
\item
  若 B 为正则环，则 A 为正则环（参见 p. 94, 习题 2 和 3）。
\item
  若 A 为正则环，且对于 A 的每个极大理想 m，B/mB 都是正则环，则 B 为正则环。（通过局部化归约到 A、B 为局部环且 ρ 为局部同态的情形。此时有 dim(B) = dim(A) + dim(B/mB)（§ 3, no. 4, 命题 7 的推论 1）。为 B 构造一个含有 dim(B) 项的正则序列。）
\end{enumerate}

\begin{enumerate}
\def\labelenumi{\arabic{enumi})}
\setcounter{enumi}{14}
\item
  设 X 为诺特空间，U 为 X 的子集。若对于 X 中每个与 U 相交的不可约闭子集 Y，U ∩ Y 都含有 Y 的一个非空开子集，则 U 在 X 中是开的。（反设 U 不是开集，考虑满足 Z ∩ U 不是开集的 X 的极小闭集 Z ⊂ X。）
\item
  设 A 为诺特环。Spec(A) 的奇异轨迹记为 Sing(A)，它是满足 \(p \in \text{Spec}(A)\) 且 \(A_p\) 不是正则环的 p 的集合。置
\end{enumerate}

\[
\operatorname{Reg} (\mathrm{A}) = \operatorname{Spec} (\mathrm{A}) - \operatorname{Sing} (\mathrm{A}).
\]

证明下列条件等价：

\begin{enumerate}
\def\labelenumi{(\roman{enumi})}
\item
  对 A 的每个商环 B，Reg(B) 在 Spec(B) 中是开集。
\item
  对 A 的每个整环商环 B，Reg(B) 都是 Spec(B) 的非空开子集。
\item
  对 A 的每个整环商环 B，Reg(B) 都含有 Spec(B) 的一个非空开子集。
\end{enumerate}

（为证明 (i) \(\Rightarrow\) (ii)，注意 \((0)\in\mathrm{Reg}(\mathbf{B})\)。为证明 (iii) \(\Rightarrow\) (i)，证明对于每个 \(\mathfrak{p}\in\mathrm{Spec}(\mathbf{A})\)，使得 \(\mathrm{V}(\mathfrak{p})\cap\mathrm{Reg}(\mathbf{A})\neq\emptyset\)，环 \(\mathbf{A}_{\mathfrak{p}}\) 是正则环（p.94, 习题 5）。使用 (iii) 并构造完全相交序列，证明 \(\mathrm{V}(\mathfrak{p})\cap\mathrm{Reg}(\mathbf{A})\) 含有 \(\mathrm{V}(\mathfrak{p})\) 的一个非空开子集。使用习题 15 得出结论。）

\begin{enumerate}
\def\labelenumi{\arabic{enumi})}
\setcounter{enumi}{16}
\tightlist
\item
  设 \(\rho: A \to B\) 为整诺特环的单同态，使 \(B\) 成为有限型 \(A\)-模。
\end{enumerate}

\begin{enumerate}
\def\labelenumi{\alph{enumi})}
\item
  假设 \(\operatorname{Reg}(\mathbf{B})\) 含有 \(\operatorname{Spec}(\mathbf{B})\) 的一个非空开子集。证明 \(\operatorname{Reg}(\mathbf{A})\) 含有 \(\operatorname{Spec}(\mathbf{A})\) 的一个非空开子集（使用 II, § 5, no. 1, 命题 2 的推论，归约到 \(\mathbf{B}\) 在 \(\mathbf{A}\) 上平坦的情形；然后在 \(\mathbf{B}\) 为正则环的情形，对适当的元素 \(x \in \mathbf{A}\) 局部化。使用习题 14, a) 得出结论。）
\item
  假设 \(\operatorname{Reg}(\mathbf{A})\) 含有 \(\operatorname{Spec}(\mathbf{A})\) 的一个非空开子集，且 \(\mathbf{K}'\)（\(\mathbf{B}\) 的分式域）是 \(\mathbf{K}\)（\(\mathbf{A}\) 的分式域）的可分扩张。证明 \(\operatorname{Reg}(\mathbf{B})\) 含有 \(\operatorname{Spec}(\mathbf{B})\) 的一个非空开子集。（取 \(\mathbf{Z} \in \mathbf{K}'\) 使得 \(\mathbf{K}' = \mathbf{K}(\mathbf{Z})\)（A, V, p.39），令 \(\mathrm{P}(\mathrm{X}) \in \mathrm{K}[\mathrm{X}]\) 为 Z 的最小多项式。用 A 中适当的元素 f 将 A 替换为环 \(A_{f}\)，证明可以归约到证明 Spec(B) 在下列情形中具有所要求的性质：
\end{enumerate}

\(\alpha)\) \(\mathbf{Z}\in \mathbf{B}\) 且 A 为正则环；

\(\beta)\) 有 \(\mathrm{P}(\mathrm{X})\in \mathrm{A}[\mathrm{X}]\)；

\(\gamma)\) 有 \(\mathbf{B} = \mathbf{A}[\mathbf{X}] / (\mathbf{P}(\mathbf{X}))\)（II, § 5, no. 1, 命题 2）；

δ) P'(Z) 在 B 中可逆（注意 P'(Z) ≠ 0 在 K' 中，因为 K' 在 K 上可分。然后使用 II, § 5, no. 1, 命题 3 的推论）。

  证明在此情形下，对于 A 的每个极大理想 m，B/mB 是一个平展 (A/m)-代数（A, V, p. 32），因而是正则环。使用习题 14, b) 得出结论。

18) 设 \(k\) 为一个域，A 为有限型 \(k\)-代数。证明 \(\operatorname{Reg}(A)\) 是 \(\operatorname{Spec}(A)\) 的开子集。（使用习题 16，归约到 A 为整环时证明 \(\operatorname{Reg}(A)\) 含有一个非空开子集。使用正规化引理（V, § 3, no. 1, 定理 1）和习题 17, a)，归约到 A 是 \(k[X_1, ..., X_n]\) 在 \(k(X_1, ..., X_n)\) 的有限拟伽罗瓦扩张 K 中的整闭包的情形。令 K' 为 K 中 \(k(X_1, ..., X_n)\) 的纯不可分闭包，A' 为 \(k[X_1, ..., X_n]\) 在 K' 中的整闭包。利用习题 17, b)，归约到为 A' 证明所要求的性质。仿照 V, § 3, no. 2 的定理 2 的证明，将 K' 嵌入扩张 \(k'(X_1^{q^{-1}}, ..., X_n^{q^{-1}}) = K''\) 中，其中 \(k'\) 是 \(k\) 的纯不可分扩张，并依据习题 17, a) 归约到 A′ 是 \(k[X_1, ..., X_n]\) 在 K'' 中的整闭包的情形。然后注意到 A' = \(k'[X_1^{q^{-1}}, ..., X_n^{q^{-1}}]\)（V, § 1, No. 3, 命题 13 的推论 2），且 A' 是正则环（p. 94, 习题 6）。

\begin{enumerate}
\def\labelenumi{\arabic{enumi})}
\setcounter{enumi}{18}
\item
  设 A 为域 k 上的有限型整环。满足 \(m \in \text{Specmax}(A)\) 且 \(A_{m}\) 为正则环的极大理想集合在 \(\text{Specmax}(A)\) 中开且稠密。
\item
  设 \(k\) 为一个域，\(k_{i}\)（\(i = 1, 2\)）为 \(k\) 的两个扩张；若超越次数有限，则令 \(n_{i}\)（\(i = 1, 2\)）为 \(k_{i}\) 在 \(k\) 上的超越次数，否则令其为符号 +\(\infty\)。证明有
\end{enumerate}

\[
\dim (k _ {1} \otimes_ {k} k _ {2}) = \inf (n _ {1}, n _ {2})
\]

并且 \(k_{1} \otimes_{k} k_{2}\) 是诺特环，当扩张 \(k_{i}\) 中至少有一个是有限型时。

\begin{enumerate}
\def\labelenumi{\arabic{enumi})}
\setcounter{enumi}{20}
\item
  设 I 为定向有序集，\((\mathbf{A}_{i}, \varphi_{ij})\) 为以 A 为归纳极限的环的归纳系统。假设对于每个 i，环 \(A_{i}\) 都是诺特环；对于每个对 \((i, j)\)（其中 i \textless{} j），\(\varphi_{ij}\) 使 \(A_{j}\) 成为平坦的 \(A_{i}\)-模；且 A 是诺特环。证明若 \(A_{i}\) 对所有 \(i \in I\) 都是正则环，则 A 是正则环。（可以先归约到 \(A_{i}\) 和 A 都是局部环且 \(\varphi_{ij}\) 是局部同态的情形。然后使用 p.94 习题 3 和 p.96 习题 14。）
\item
  \begin{enumerate}
  \def\labelenumii{\alph{enumii})}
  \tightlist
  \item
    设 \(k\) 为一个域，A 为有限型 \(k\)-代数，\(k'\) 为 \(k\) 的有限可分扩张。证明若 A 为正则环，则 \(A \otimes_k k'\) 为正则环。（可以使用 p.96 习题 14, b)。）
  \end{enumerate}
\end{enumerate}

\begin{enumerate}
\def\labelenumi{\alph{enumi})}
\setcounter{enumi}{1}
\tightlist
\item
  在 A 具有相同假设的条件下，设 \(k'\) 为 \(k\) 的有限纯不可分扩张。证明一般而言 \(A \otimes_k k'\) 不是正则环（取 \(A = k'\)）。
\end{enumerate}

\begin{enumerate}
\def\labelenumi{\arabic{enumi})}
\setcounter{enumi}{22}
\tightlist
\item
  设 \(k\) 为特征指数为 \(p\) 的域，A 为有限型 \(k\)-代数，\(\mathfrak{p}\) 为 A 的素理想。证明下列条件等价：
\end{enumerate}

\begin{enumerate}
\def\labelenumi{(\roman{enumi})}
\item
  环 \(A_{\mathfrak{p}} \otimes_{k} k'\) 对 k 的每个扩张 \(k'\) 都是正则环；
\item
  环 \(A_{\mathfrak{p}} \otimes_{k} k^{p^{-\infty}}\) 是正则环；
\item
  环 \(\mathbf{A}_{\mathfrak{p}} \otimes_{k} k'\) 对每个有限且纯不可分的扩张 \(k'\)（扩张 \(k\)）都是正则环。
\end{enumerate}

 （为证明 (ii) \(\Rightarrow\) (iii)，使用 p.96 习题 14, a)。为证明 (iii) \(\Rightarrow\) (i)，只需证明 \(A_{\mathfrak p} \otimes_k k'\) 在 \(k'\) 是 \(k\) 的有限生成扩张时为正则环（习题 21），并且当 \(k'\) 是纯超越扩张的纯不可分扩张时（习题 22 和 14, a)），且被纯超越扩张的一个有限拟伽罗瓦扩张所控制时也是如此。然后把 \(k'\) 控制在扩张 \(k''(\mathrm{X}_1, ..., \mathrm{X}_n)\) 中，其中 \(k''\) 是 \(k\) 的有限纯不可分扩张，并归约到 \(k' = k''(\mathrm{X}_1, ..., \mathrm{X}_n)\) 的情形。由习题 14, a)，置 \(\mathbf{B} = \mathbf{A}_{\mathfrak{p}} \otimes_k k''\)，根据 (iii) 得到一个正则代数，而 \(\mathbf{A}_{\mathfrak{p}} \otimes_k k' = \mathbf{B} \otimes_{k''} k''(\mathbf{X}_1, ..., \mathbf{X}_n)\) 根据 p.94 习题 6, e) 是正则代数。）

若 \(p \in \text{Spec}(A)\) 具有上述等价性质，则称 p 为光滑点。

\begin{enumerate}
\def\labelenumi{\arabic{enumi})}
\setcounter{enumi}{23}
\item
  设 \(k\) 为一个域，A 为 \(k\)-代数，\(k'\) 为 \(k\) 的纯不可分扩张，且 \(\mathbf{A}' = \mathbf{A} \otimes_k k'\)。证明典范映射 \(\operatorname{Spec}(\mathbf{A}') \to \operatorname{Spec}(\mathbf{A})\) 是同胚。（当 A 为域时，\(\mathbf{A}'\) 的维数为 0（习题 20）；证明它只有一个素理想。由此可推出，对于 A 的每个素理想 \(\mathfrak{p}\)，\(\mathfrak{p}'\) 是 \(\mathfrak{p}\mathbf{A}'\) 的根基，且映射 \(\mathfrak{p} \mapsto \mathfrak{p}'\) 是连续的。）
\item
  设 \(k\) 为一个域，A 为有限型 \(k\)-代数，\(\operatorname{Lis}(A)\) 为由光滑的 \(\mathfrak{p} \in \operatorname{Spec}(A)\) 构成的集合（习题 23）。证明 \(\operatorname{Lis}(A)\) 在 \(\operatorname{Spec}(A)\) 中是开集。（使用习题 18、23 和 24。）
\item
  设 \(k\) 为一个域，A 为有限型整环 \(k\)-代数。证明 Lis(A) 稠密当且仅当 A 是可分的 \(k\)-代数；等价地，A 的分式域是 \(k\) 的可分扩张（A, V, p.115）。
\item
  设 \(k\) 为一个域，A 为有限型且整的可分 \(k\)-代数。A 的光滑极大理想集合是 Specmax(A) 的开稠密子集。
\item
  设 \(k\) 为一个域，\(k'\) 为 \(k\) 的扩张。证明对于每个扩张 \(k''\)（它是 \(k\) 的扩张），\(k' \otimes_k k''\) 是正则诺特环当且仅当 \(k'\) 是 \(k\) 的有限型可分扩张。
\end{enumerate}

29) 设 k 为一个域，A 为完备诺特局部 k-代数，\(\kappa_{A}\) 为 A 的剩余域。

\begin{enumerate}
\def\labelenumi{\alph{enumi})}
\item
  假设 \(\kappa_{\mathbf{A}}\) 是 \(k\) 的可分扩张。证明存在 \(k\)-代数同态，从 \(\kappa_{\mathbf{A}}\) 到 A，它是从 A 到 \(\kappa_{\mathbf{A}}\) 的典范同态的截面（IX, § 3, n° 2, 命题 1）。
\item
  假设 \(\kappa_{\mathbf{A}}\) 是 \(k\) 的纯不可分扩张。证明不一定存在 \(k\)-代数截面，从 \(\kappa_{\mathbf{A}}\) 到 A。（取特征为 2 的域 \(k\)，取不是平方的元素 \(a\)，它属于 \(k\)（例如 \(k = \mathbf{F}_2(\mathrm{T})\)，\(a = \mathrm{T}\)），取理想 \(\mathfrak{p}\)，它是 \(k[\mathrm{X}]\) 中由 \(\mathrm{X}^2 + a\) 生成的理想，并置 \(\mathrm{A} = \widehat{k[\mathrm{X}]_{\mathfrak{p}}}\)。）
\item
  假设 A 为正则环，且 \(\kappa_{\mathbf{A}}\) 是 \(k\) 的可分扩张。证明 A 作为 \(k\)-代数同构于 \(\kappa_{\mathbf{A}}[[\mathrm{T}_{1}, ..., \mathrm{T}_{n}]]\)，其中 \(n = \dim(\mathbf{A})\)。
\item
  证明当 \(\kappa_{\mathrm{A}}\) 不是 \(k\) 的可分扩张时，情形不一定如此。
\end{enumerate}

\begin{enumerate}
\def\labelenumi{\arabic{enumi})}
\setcounter{enumi}{29}
\item
  设 \(k\) 为一个域，A 为局部、诺特且完备的 \(k\)-代数。若对于每个有限扩张 \(k'\)（它是 \(k\) 的扩张），环 \(A \otimes_k k'\) 都是正则环，则称 A 形式光滑。证明 A 形式光滑当且仅当对于每个有限纯不可分扩张 \(k'\)（它是 \(k\) 的扩张），环 \(A \otimes_k k'\) 都是正则环（可以使用 p.96 习题 14，并仿照 p.97 习题 23 的证明）。
\item
  设 \(k\) 为一个域。
\end{enumerate}

\begin{enumerate}
\def\labelenumi{\alph{enumi})}
\item
  设 A 为诺特、局部、完备且正则的 \(k\)-代数，其剩余域 \(\kappa_{\mathrm{A}}\) 是 \(k\) 的可分扩张。证明 A 形式光滑且同构于 \(\kappa_{\mathrm{A}}[[\mathbf{T}_{1}, ..., \mathbf{T}_{n}]]\)。
\item
  设 B 为有限型 \(k\)-代数，\(\mathfrak{p} \in \operatorname{Spec}(\mathbf{B})\) 为光滑点（p.97, 习题 23）。证明 \(\hat{\mathbf{B}}_{\mathfrak{p}}\) 是形式光滑的 \(k\)-代数。
\item
  设 \(\kappa\) 为 \(k\) 的有限生成扩张，且 \(n\in\mathbf{N}\)。证明存在形式光滑的 \(k\)-代数，其剩余域为 \(\kappa\)，维数为 \(n\)。
\item
  设 A 为形式光滑的 k-代数，其剩余域 \(\kappa_{A}\) 不是 k 的可分扩张。证明 A 作为 k-代数不同构于 \(\kappa_{A}[[T_{1},...,T_{n}]]\)。
\end{enumerate}

32) 设 \(k\) 为一个域，K 为 \(k\) 的有限生成扩张，\(n\) 为正整数，A、B 为两个形式光滑 \(k\)-代数，维数为 \(n\)，且 \(\kappa_{A}\) 和 \(\kappa_{B}\) 作为 \(k\) 的扩张都同构于 K。拟证明 A、B 是同构的 \(k\)-代数。

\begin{enumerate}
\def\labelenumi{\alph{enumi})}
\item
  归约到 K 是 k 的有限纯不可分扩张的情形（习题 29, a)）。
\item
  考虑 \(C = A \hat{\otimes}_{k} B\)，它是 \(A \otimes_{k} B\) 关于理想的完备化代数
\end{enumerate}

\[
\left(\mathbf {A} \otimes_ {k} \mathfrak {m} _ {\mathbf {B}}\right) + \left(\mathfrak {m} _ {\mathbf {A}} \otimes_ {k} \mathbf {B}\right).
\]

证明 C 是诺特、局部、完备且正则的 \(k\)-代数。（为证明最后一点，使用从 A 到 C 的典范同态 \(\rho\)。我们将证明它使 C 成为平坦的 A-模，且 \(C \otimes_{\mathbf{A}} \kappa_{\mathbf{A}}\) 同构于 \(\kappa_{\mathbf{A}} \otimes_{k} \mathbf{B}\)。然后使用 p.96 习题 14。）

\begin{enumerate}
\def\labelenumi{\alph{enumi})}
\setcounter{enumi}{2}
\tightlist
\item
  证明 \(\rho\) 在剩余域上诱导同构，并将 \(m_{\mathbf{A}}/m_{\mathbf{A}}^{2}\) 单射到 \(m_{\mathbf{C}}/m_{\mathbf{C}}^{2}\) 中。然后构造 \(\pi\) 作为 \(\rho\) 的一个回缩，它是 \(k\)-代数同态。证明 \(\pi \circ \rho': B \to A\) 是 \(k\)-代数同构。（我们将从 B 到 C 的典范同态记为 \(\rho'\)。）
\end{enumerate}

\exercisesection{分次环的维数}

\begin{enumerate}
\def\labelenumi{\arabic{enumi})}
\tightlist
\item
  设 A 为一个环，p 为 A{[}T{]} 的齐次素理想，且 \(p_{0}\) 为 A 中的素理想 \(p \cap A\)。
\end{enumerate}

\begin{enumerate}
\def\labelenumi{\alph{enumi})}
\item
  若 \(T \in p\)，则 \(p = p_{0} + T \cdot A[T]\)。
\item
  若 T \(\notin\) p，则 \(\mathfrak{p} = \mathfrak{p}_0\)。A{[}T{]}。
\item
  推出在情形 \(a)\) 下有 \(\operatorname{htgr}(\mathfrak{p}) = \operatorname{ht}(\mathfrak{p}_0) + 1\)，而在情形 \(b)\) 下有 \(\operatorname{htgr}(\mathfrak{p}) = \operatorname{ht}(\mathfrak{p}_0)\)。
\item
  证明 \(\dim \operatorname{gr}(\mathbf{A}[\mathbf{T}]) = \dim (\mathbf{A}) + 1\)。
\item
  由此推出一个满足 dim(B) ≠ dimgr(B) 的分次环 B 的例子（p.84, 习题 7）。
\end{enumerate}

\begin{enumerate}
\def\labelenumi{\arabic{enumi})}
\setcounter{enumi}{1}
\tightlist
\item
  设 A 为一个环，\(\mathcal{F} = (\mathcal{F}_i)_{i \in \mathbf{Z}}\) 为 A 的递减滤过，且 \(\mathcal{F}_i = \mathbf{A}\) 对 \(i \leqslant 0\) 成立。令 \(\mathcal{A}\) 为 \(\bigoplus_{i \in \mathbf{Z}} \mathcal{F}_i \cdot v^{-i}\) 在 \(\mathbf{A}[v, v^{-1}]\) 中的子环，令 \(\mathrm{gr}_{\mathcal{F}}(\mathbf{A})\) 为
\end{enumerate}

该滤过环 (A, F) 的伴随分次环。

\begin{enumerate}
\def\labelenumi{\alph{enumi})}
\item
  证明同态 \(f: \mathcal{A} \to \mathrm{gr}_{\mathcal{F}}(\mathbf{A})\) 由 \(f\left(\sum_{i} a_i v^{-i}\right) = \sum_{i} \overline{a}_i\) 定义，其中 \(\overline{a}_i\) 是 \(a_i\) 在 \(\mathcal{F}_i / \mathcal{F}_{i+1}\) 中的类，并诱导出 \(\mathcal{A}/v. \mathcal{A}\) 与 \(\mathrm{gr}_{\mathcal{F}}(\mathbf{A})\) 之间的同构。
\item
  证明对于 A 的每个可逆元素 \(v_{0}\)，由下式定义的同态 \(e_{v_{0}}: \mathcal{A} \to \mathbf{A}\)
\end{enumerate}

\[
e _ {v _ {0}} (\sum a _ {i} v ^ {- i}) = \sum a _ {i} v _ {0} ^ {- i}
\]

诱导出 \(\mathcal{A}/(v - v_{0})\mathcal{A}\) 与 A 之间的同构。

\begin{enumerate}
\def\labelenumi{\alph{enumi})}
\setcounter{enumi}{2}
\item
  证明 \(\mathcal{A}/\mathrm{A}[v]\) 是挠 \(\mathrm{A}[v]\)-模。
\item
  现在假设 A 含有域 k，证明复合态射
\end{enumerate}

\[
k [ v ] \rightarrow \mathrm{A} [ v ] \rightarrow \mathcal {A} \quad (\text {where } \mathrm{A} [ v ] = \bigoplus_ {i \leqslant 0} \mathcal {F} _ {i}. v ^ {- i})
\]

使 \(\mathcal{A}\) 成为无挠的 k{[}v{]}-模。推出 \(\mathcal{A}\) 是平坦的 k{[}v{]}-模。

\begin{enumerate}
\def\labelenumi{\alph{enumi})}
\setcounter{enumi}{4}
\tightlist
\item
  假设 A 是局部环，极大理想为 m，并含有域 k，使得 A（作为加法群）是 k 与 m 的直和。证明 A 存在理想 S，使得 \(\mathcal{A}\) 是 k{[}v{]} 与 S 的直和（我们置 \(\mathcal{F}_{i} = m^{i}\) 对每个 \(i \in \mathbf{Z}\) 成立）。
\end{enumerate}

\begin{enumerate}
\def\labelenumi{\arabic{enumi})}
\setcounter{enumi}{2}
\tightlist
\item
  继续使用上一习题的记号。设 M 为一个 A-模，带有滤过 \((\mathbf{K}_i)_{i\in\mathbf{Z}}\)，使得 \(\mathbf{K}_i = \mathbf{M}\) 对 \(i \leqslant 0\) 成立。
\end{enumerate}

\begin{enumerate}
\def\labelenumi{\alph{enumi})}
\tightlist
\item
  证明分次 A-模 \(\mathcal{M} = \sum_{i\in \mathbf{Z}}\mathrm{K}_{i}v^{-i}\) 具有自然的 \(\mathcal{A}\)-分次模结构。
\end{enumerate}



\begin{enumerate}
\def\labelenumi{\alph{enumi})}
\setcounter{enumi}{1}
\item
  证明分次 A-代数 \(\mathcal{A}\) 由其次数为 1 和 -1 的元素生成，当且仅当 A 存在理想 q 使得 \(\mathcal{F}_i = \mathfrak{q}^i\)（\(i \in \mathbf{Z}\)）。
\item
  假设最后一个条件成立，证明滤过 \((\mathbf{K}_i)_{i\in \mathbb{Z}}\) 是 q-良好的，当且仅当 \(\mathcal{A}\)-分次模 \(\mathcal{M}\) 是有限型的。
\item
  证明若滤过 \((\mathbf{K}_{i})_{i\in\mathbf{Z}}\) 是 q-良好的，则 A-模 M 是 v-无挠的；反之，给定一个无 v-挠的有限生成 A-模 M，可以在其上赋予一个 q-良好滤过，使得伴随 A-模为 M。
\end{enumerate}

\begin{enumerate}
\def\labelenumi{\arabic{enumi})}
\setcounter{enumi}{3}
\tightlist
\item
  采用习题 2 的记号和假设。称滤过 \(\mathcal{F}\) 为准幂滤过，若存在整数 \(k > 0\) 使得
\end{enumerate}

\[
\mathcal {F} _ {n} = \sum_ {i = 1} ^ {k} \mathcal {F} _ {i}. \mathcal {F} _ {n - i} \quad \text { for } \quad n \geqslant 1.
\]

证明若环 A 是诺特环，则下列条件等价：

\begin{enumerate}
\def\labelenumi{(\roman{enumi})}
\item
  滤过 \(\mathcal{F}\) 是准幂滤过。
\item
  环 \(\mathcal{A}\) 是诺特环。
\item
  理想 \(\mathcal{N} = \sum_{i > 1} v^{-i}\mathcal{F}_i\) 是 \(\mathcal{A}\) 的有限型理想。
\item
  存在整数 \(k \geqslant 1\)，使得 \(\mathcal{F}\) 是 A 的关于序关系“\(\mathcal{H}_n \subset \mathfrak{G}_n\) 对每个 \(n\) 都成立”的、以 \(\mathcal{F}_1, ..., \mathcal{F}_k\) 为首项的最小滤过。
\item
  A-代数 \(\mathcal{A}\) 是有限型的。
\end{enumerate}

\begin{enumerate}
\def\labelenumi{\arabic{enumi})}
\setcounter{enumi}{4}
\tightlist
\item
  设 \(\rho: A \to B\) 为诺特环的单同态，使 \(B\) 成为有限型 \(A\)-代数。称 \(B\) 在 \(A\) 上的超越次数为 \(d_A(B)\)，其中 \(d\) 是满足存在 \(A\)-代数单同态 \(b: A[X_1, \ldots, X_d] \to B\) 的最大整数。a) 证明若 \(\sigma: B \to C\) 满足与 \(\rho\) 相同的假设，则有不等式
\end{enumerate}

\[
\mathrm{d} _ {\mathbf {A}} (\mathbf {C}) \geqslant \mathrm{d} _ {\mathbf {A}} (\mathbf {B}) + \mathrm{d} _ {\mathbf {B}} (\mathbf {C}).
\]

\begin{enumerate}
\def\labelenumi{\alph{enumi})}
\setcounter{enumi}{1}
\item
  证明若 A、B 是整环，则 \(d_{\mathbf{A}}(\mathbf{B})\) 等于 B 的分式域在 A 的分式域上的超越次数。此外，若在 a) 中假设 C 是整环，则有等式 \(d_{\mathbf{A}}(\mathbf{C}) = d_{\mathbf{A}}(\mathbf{B}) + d_{\mathbf{B}}(\mathbf{C})\)。
\item
  证明若 B 是 A{[}v, \(v^{-1}\) {]} 的一个包含 A{[}v{]} 的子代数，其中 v 是不定元，则有 \(d_{\mathbf{A}}(\mathbf{B}) = 1\)。
\end{enumerate}

  6) 设 \(\rho: A \to B\) 为诺特环的单同态，使 \(B\) 成为有限型 \(A\)-代数。设 \(q\) 为 \(B\) 中异于 \(B\) 的素理想，并置 \(p = \rho^{-1}(q)\)。a) 证明使用上一习题的记号，有不等式

\[
(*) \quad \mathrm{ht} (\mathfrak {q}) + \mathrm{d} _ {\mathrm{A} / \mathfrak {p}} (\mathrm{B} / \mathfrak {q}) \leqslant \mathrm{ht} (\mathfrak {p}) + \mathrm{d} _ {\mathrm{A}} (\mathrm{B}).
\]

（可以对 \(d_{A}(B)\) 使用归纳法。）

\begin{enumerate}
\def\labelenumi{\alph{enumi})}
\setcounter{enumi}{1}
\item
  证明若 A-代数 B 由 d 个元素生成且 \(\mathrm{d}_{\mathbf{A}}(\mathbf{B}) = d\)，则 (*) 中为等式。
\item
   证明若 A 是整环，且对于每个 n，多项式环 A{[}T \(_{1}\) , \ldots, T \(_{n}\) {]} 都是串联环，则对于每个有限型整 A-代数 B 及 B 中异于 B 的每个素理想 q，(*) 中为等式。
\end{enumerate}

\begin{enumerate}
\def\labelenumi{\arabic{enumi})}
\setcounter{enumi}{6}
\tightlist
\item
  设 A 为诺特环，\(\mathcal{F} = (\mathcal{F}_i)_{i\in \mathbf{Z}}\) 为 A 的递减滤过，使得 \(\mathcal{F}_0 = \mathbf{A}\)、\(\mathcal{F}_1 \neq \mathbf{A}\)，且 A-代数 \(\mathbf{B} = \sum_{i=1}^{\infty} \mathcal{F}_i v^{-i}\) 是有限型的。
\end{enumerate}

我们拟证明 \(\dim(B) = \dim(A) + 1\)，并利用这一事实推广 § 6, no. 3, 命题 5 的推论。将使用前面两个习题的结果。

\begin{enumerate}
\def\labelenumi{\alph{enumi})}
\tightlist
\item
  设 M 为 B 的极大理想。证明有
\end{enumerate}

\[
\operatorname{ht} (\mathfrak {M}) \leqslant \dim (\mathrm{A}) + 1,
\]

并由此推出不等式

\[
\dim (\mathbf {B}) \leqslant \dim (\mathbf {A}) + 1.
\]

\begin{enumerate}
\def\labelenumi{\alph{enumi})}
\setcounter{enumi}{1}
\tightlist
\item
  利用 A{[}v{]}-模 A{[}v, \(v^{-1}\) {]}/A{[}v{]} 的每个元素都被 v 的某个幂消去这一事实，证明有
\end{enumerate}

\[
\dim (\mathrm{A} [ v, v ^ {- 1} ]) = \dim (\mathrm{A} [ v ]) = \dim (\mathrm{A}) + 1,
\]

并推出 \(\dim(\mathbf{B}) \geqslant \dim(\mathbf{A}) + 1\)，从而最终有 \(\dim(\mathbf{B}) = \dim(\mathbf{A}) + 1\)（可以使用环 \(\mathbf{A}[v, v^{-1}]\) 是 \(\mathbf{B}\) 的分式环这一事实）。c) 证明 \(\mathbf{B}\) 的极小素理想是如下形式的理想：

\[
\sum_ {i \in \mathbf {Z}} \left(\mathfrak {p} \cap \mathcal {F} _ {i}\right) v ^ {- i}
\]

其中 p 是 A 的极小素理想。

\begin{enumerate}
\def\labelenumi{\alph{enumi})}
\setcounter{enumi}{3}
\item
  证明 \(v\mathbf{B} \neq \mathbf{B}\)，且 \(v\) 不属于 B 的任何极小素理想。推出不等式 \(\dim(\mathbf{B}/v\mathbf{B}) \leqslant \dim(\mathbf{A})\)（可以使用 § 1 及 § 3, no. 3, 命题 4 的推论 1）。
\item
  现在假设 A 存在包含 \(\mathcal{F}_1\) 的极大理想 m，且 ht(m) = dim(A)。证明在此情形下有 dim(B/vB) ≥ dim(A)，因而 dim(B/vB) = dim(A)。（可以归约到 A 为极大理想 m 的局部环的情形，证明由 v 和 \(\sum_{i\in\mathbf{Z}}(\mathcal{F}_i\cap m)v^{-i}\) 生成的 B 的理想 M 是高度为 dim(A) + 1 的 B 的极大理想，并在 M 处局部化 B。）
\end{enumerate}





\begin{enumerate}
\def\labelenumi{\alph{enumi})}
\setcounter{enumi}{5}
\tightlist
\item
  若 A 为诺特局部环，且 F 为 A 的递减滤过，使得 \(F_{0} = A\)、\(F_{1} \neq A\)，则有等式
\end{enumerate}

\[
\dim (\operatorname{gr} (A)) = \dim (A).
\]

\begin{enumerate}
\def\labelenumi{\arabic{enumi})}
\setcounter{enumi}{7}
\tightlist
\item
  设 A 为诺特环，且 B 为多项式环 A{[}T{]} 的分次子-A-代数。
\end{enumerate}

\begin{enumerate}
\def\labelenumi{\alph{enumi})}
\tightlist
\item
  证明有不等式
\end{enumerate}

\[
\dim (\mathbf {A}) \leqslant \dim (\mathbf {B}) \leqslant \dim (\mathbf {A}) + 1.
\]

\begin{enumerate}
\def\labelenumi{\alph{enumi})}
\setcounter{enumi}{1}
\tightlist
\item
  设 \(B_{i}\) 为 \(b \in A\) 中满足 \(bT^{i} \in B\) 的元素的集合。证明若有 \(\dim(B) = \dim(A)\)，则存在正整数 k，使得对于所有 \(i \geqslant k\)，\(B_{i}\) 都包含于 A 的极小素理想的交中。
\end{enumerate}

\begin{enumerate}
\def\labelenumi{\arabic{enumi})}
\setcounter{enumi}{8}
\item
  设 A 为一个环，M 为有限表示的 A-模。对于每个 \(\mathfrak{p} \in \operatorname{Spec}(A)\)，置 \(\kappa(\mathfrak{p}) = A_{\mathfrak{p}} / \mathfrak{p} \cdot A_{\mathfrak{p}}\)。证明对于所有 \(k \in \mathbf{N}\)，集合 \(\mathfrak{p} \in \operatorname{Spec}(A)\) 中满足 \(\dim_{\kappa(\mathfrak{p})} (M \otimes_A \kappa(\mathfrak{p})) \geqslant k\) 的元素构成 \(\operatorname{Spec}(A)\) 的闭子集。（可以取 M 的表示 \(\mathrm{A}^m \xrightarrow{u} \mathrm{A}^n \longrightarrow \mathrm{M} \longrightarrow 0\)，并考察矩阵 \(u\) 的子式。）
\item
  设 \(\rho: A \to B\) 为环同态。对于每个 \(\mathfrak{p} \in \operatorname{Spec}(B)\)，置
\end{enumerate}

\[
\mathrm{d} (\mathfrak {p}) = \dim_ {\mathfrak {p}} \left(\mathbf {B} \otimes_ {\mathbf {A}} \kappa \left(\rho^ {- 1} (\mathfrak {p})\right)\right).
\]

  本习题旨在证明，当 B 是有限生成的 A-代数时，映射 \(p \mapsto d(p)\) 从 \(\operatorname{Spec}(B)\) 到 \(\mathbf{R}\) 是上半连续的（“Chevalley 定理”）。

记 \(\delta \subset \mathbf{B} \otimes_{\mathbf{A}} \mathbf{B}\) 为典范 A-同态 \(\mathbf{B} \otimes_{\mathbf{A}} \mathbf{B} \to \mathbf{B}\) 的核。对于每个整数 \(r\)，置

\[
\mathrm{P} _ {\mathbf {B} / \mathbf {A}} ^ {r} = (\mathbf {B} \otimes_ {\mathbf {A}} \mathbf {B}) / \delta^ {r}.
\]

称 B ⊗\_A B-代数 P \(_{B/A}^{r}\) 为 r 阶主部代数。通过同态 b ↦ b ⊗ 1 将 B ⊗\_A B 看作 B-代数，于是对于每个 r，P \(_{B/A}^{r}\) 是 B-代数，且典范商同态 P \(_{B/A}^{r+1}\) → P \(_{B/A}^{r}\) 定义出 B-代数的射影系统。

\(^{1}\) 这是作为 \(\kappa(\mathfrak{p})\) 上向量空间的维数。

\begin{enumerate}
\def\labelenumi{\alph{enumi})}
\tightlist
\item
  设 \(p: \mathbf{A} \to \mathbf{A}'\) 为环同态。置 \(\mathbf{B}' = \mathbf{B} \otimes_{\mathbf{A}} \mathbf{A}'\)，并以 \(q: \mathbf{B} \to \mathbf{B}'\) 表示典范同态。设 \(r\) 为整数。证明唯一的 \(\mathbf{B}'\)-代数同态
\end{enumerate}

\[
\mathbf {P} _ {\mathbf {B} / \mathbf {A}} ^ {r} \otimes_ {\mathbf {B}} \mathbf {B} ^ {\prime} \rightarrow \mathbf {P} _ {\mathbf {B} ^ {\prime} / \mathbf {A} ^ {\prime}} ^ {r}
\]

它将 \((b_{1} \otimes b_{2}) \otimes 1\) 的类（\(b_{i} \in \mathbf{B}\)）对应到 \(q(b_{1}) \otimes q(b_{2})\) 的类，是一个同构。b) 设 S 为 B 的一个乘法子集。证明典范同态

\[
\mathrm{S} ^ {- 1} \mathrm{P} _ {\mathrm{B/A}} ^ {r} \rightarrow \mathrm{P} _ {\mathrm{S} ^ {- 1} \mathrm{B/A}} ^ {r}
\]

是一个同构。

\begin{enumerate}
\def\labelenumi{\alph{enumi})}
\setcounter{enumi}{2}
\item
  假设 B 是有限表示的 A-代数，即同构于 \(\mathrm{A}[\mathrm{T}_{1},\ldots,\mathrm{T}_{n}]\) 对某个有限生成理想的商。证明对于每个 \(r\in\mathbf{N}\)，\(\mathbf{P}_{\mathrm{B}/\mathrm{A}}^{r}\) 是有限表示的 B-模。对于每个 \(\mathfrak{p}\in\operatorname{Spec}(\mathbf{B})\)，置 \(\mathrm{P}_{\mathrm{B}/\mathrm{A}}^{\infty}[\mathfrak{p}]=\lim_{r}\mathrm{P}_{\mathrm{B}/\mathrm{A}}^{r}\otimes_{\mathrm{B}}\kappa(\mathfrak{p})\)。证明 \(\mathrm{P}_{\mathrm{B}/\mathrm{A}}^{\infty}[\mathfrak{p}]\) 是一个 \(\kappa(\mathfrak{p})\)-代数，它是局部且完备的，剩余域为 \(\kappa(\mathfrak{p})\)。令 \(k\in\mathbf{N}\)。证明集合 \(\mathfrak{p}\in\operatorname{Spec}(\mathbf{B})\) 中满足 \(\dim(\mathrm{P}_{\mathrm{B}/\mathrm{A}}^{\infty}[\mathfrak{p}])\geqslant k\) 的元素构成 \(\operatorname{Spec}(\mathbf{B})\) 的闭子集。（对于每个 \(r\in\mathbf{N}\)，令 \(\mathrm{F}_{r,k}\) 为 \(\mathfrak{p}\) 中满足 \(\dim_{\kappa(\mathfrak{p})}(\mathrm{P}_{\mathrm{B}/\mathrm{A}}^{r}\otimes\kappa(\mathfrak{p}))\geqslant\binom{r+k-1}{r-1}\) 的元素集合。我们将证明 \(\mathrm{F}_{r,k}\) 是闭集（习题 9），从而 \(\Phi_{k}=\bigcap_{r}\mathrm{F}_{r,k}\) 是闭集。然后应用 § 6, no. 3, 命题 6 的推论 2。）
\item
  设 m 为 B 的一个理想，\(\mu\subset(\mathrm{B/m})\otimes_{\mathrm{A}}\mathrm{B}\) 为典范映射 \((\mathrm{B/m})\otimes_{\mathrm{A}}\mathrm{B}\to\mathrm{B/m}\) 的核。证明有典范同构
\end{enumerate}

\[
\mathrm{P} _ {\mathrm{B/A}} ^ {r} \otimes_ {\mathrm{B}} (\mathrm{B/m}) \rightarrow ((\mathrm{B/m}) \otimes_ {\mathrm{A}} \mathrm{B}) / \mu^ {r}.
\]

假设 A 为域，B 为有限生成的 A-代数，且 m 为 B 的极大理想。证明 \(P_{B/A}^{\infty}[m]\) 同构于 \((B/m)\otimes_{A}B\) 在极大理想 \(\mu\) 处的完备局部化。推出 \(\dim(P_{B/A}^{\infty}[m])=\dim_{m}(B)\)。（应用零点定理和 § 2, no. 3 的定理 1。注意 A 的扩张 B/m 包含于 A 的一个有限纯不可分扩张的有限伽罗瓦扩张中。）

\begin{enumerate}
\def\labelenumi{\alph{enumi})}
\setcounter{enumi}{4}
\tightlist
\item
  假设 B 为有限生成的 A-代数。设 \(p \in \operatorname{Spec}(B)\)。证明有
\end{enumerate}

\[
\dim (\mathrm{P} _ {\mathrm{B/A}} ^ {\infty} [ \mathfrak {p} ]) = \dim \left(\mathrm{B} \otimes_ {\mathrm{A}} \kappa (\rho^ {- 1} (\mathfrak {p}))\right).
\]

  首先注意对于每个 \(r \in \mathbf{N}\)，有 \(\mathbf{P}_{\mathbf{B}/\mathbf{A}}^{\mathbf{r}} \otimes_{\mathbf{B}} \kappa(\mathfrak{p}) = (\mathbf{P}_{\mathbf{B}/\mathbf{A}}^{\mathbf{r}} \otimes_{\mathbf{B}} \mathbf{B}') \otimes_{\mathbf{B}'} \kappa(\mathfrak{p})\)，这是通过置 \(\mathbf{B}' = \mathbf{B} \otimes_{\mathbf{A}} \kappa(\rho^{-1}(\mathfrak{p}))\) 得到的，因而根据 \(a\) )，有 \(\mathbf{P}_{\mathbf{B}/\mathbf{A}}^{\infty}[\mathfrak{p}] = \mathbf{P}_{\mathbf{B}' / \kappa(\rho^{-1}(\mathfrak{p}))}^{\infty}[\mathfrak{p}\mathbf{B}']\)。因此归约到 A 为域时证明等式

\[
\dim \left(P _ {B / A} ^ {\infty} [ p ]\right) = \dim_ {p} (B).
\]

注意当 p 为极大理想时，此等式由 d) 得出。一般地，设 \(q_{1}, \ldots, q_{l}\) 为包含于 p 中的 B 的极小素理想。使用 Nullstellensatz 证明，存在 \(V(p)\) 的稠密开集 U，使得对于属于 U 的每个极大理想 m，包含于 m 中的极小素理想都包含于 U。使用 c) 证明，缩小 U 后还可以假设对于每个 \(p' \in U\)，有 \(\dim(\mathrm{P}_{\mathrm{B/A}}^{\infty}[\mathfrak{p}]) = \dim(\mathrm{P}_{\mathrm{B/A}}^{\infty}[\mathfrak{p}'])\)。得出结论。）

\begin{enumerate}
\def\labelenumi{\alph{enumi})}
\setcounter{enumi}{5}
\tightlist
\item
  假设 B 为有限生成的 A-代数。证明函数 \(\mathfrak{p} \mapsto \mathrm{d}(\mathfrak{p}) = \dim_{\mathfrak{p}} (\mathrm{B} \otimes_{\mathbf{A}} \kappa(\rho^{-1}(\mathfrak{p})))\) 是上半连续的。（当 B 有限表示时，使用 \(e\) 和 \(c\)。一般地，B 具有形式 \(\mathrm{A}[\mathrm{T}_1, ..., \mathrm{T}_n]/\mathfrak{J}\)。令 \(\mathfrak{J}_{\alpha}\) 为 \(\mathrm{A}[\mathrm{T}_1, ..., \mathrm{T}_n]\) 中包含于 \(\mathfrak{J}\) 的有限生成理想族，并按包含关系排序，置 \(\mathrm{B}_{\alpha} = \mathrm{A}[\mathrm{T}_1, ..., \mathrm{T}_n]/\mathfrak{J}_{\alpha}\)。对于每个 \(\alpha\)，有 \(\operatorname{Spec}(\mathbf{B}) \subset \operatorname{Spec}(\mathrm{B}_{\alpha}) \subset \operatorname{Spec}(\mathrm{A}[\mathrm{T}_1, ..., \mathrm{T}_n])\)，且 \(\operatorname{Spec}(\mathbf{B}) = \bigcap \operatorname{Spec}(\mathrm{B}_{\alpha})\)。对于每个 \(\mathfrak{p} \in \operatorname{Spec}(\mathbf{B})\)，置
\end{enumerate}

\[
\mathrm{d} _ {\alpha} (\mathfrak {p}) = \dim_ {\mathfrak {p}} \left(\mathrm{B} _ {\alpha} \otimes \kappa \left(\rho^ {- 1} (\mathfrak {p})\right)\right).
\]

  证明有 \(d_{\alpha}(p) \geqslant d(p)\) 及 \(d(p) = \inf_{\alpha} d_{\alpha}(p)\)，这是根据 \(\text{Spec}(\kappa(\rho^{-1}(p))[T_1, \ldots, T_n])\) 的诺特性。推出 \(p \mapsto d(p)\) 是一个递减滤族上半连续函数的下确界。）

\begin{enumerate}
\def\labelenumi{\arabic{enumi})}
\setcounter{enumi}{10}
\tightlist
\item
  设 \(\rho: \mathbf{A} \to \mathbf{B}\) 为环同态，\(^a\rho: \operatorname{Spec}(\mathbf{B}) \to \operatorname{Spec}(\mathbf{A})\) 为相应映射。
\end{enumerate}

\begin{enumerate}
\def\labelenumi{\alph{enumi})}
\item
  证明对于每个 \(\mathfrak{q} \in \operatorname{Spec}(\mathbf{A})\)，\(\operatorname{Spec}(\mathbf{B} \otimes_{\mathbf{A}} \kappa(\mathfrak{q}))\) 可与 \(^a\rho^{-1}(\mathfrak{q})\) 识别。
\item
  假设 B 为有限生成的 A-代数。证明 p 是其纤维 \(^a\rho^{-1}(^a\rho(p))\) 中的孤立点，当且仅当 \(\dim_p(B \otimes_A \kappa(^a\rho(p)))\) 为零。
\item
  从习题 10 推出，在 \(b\) 的假设下，\(\operatorname{Spec}(\mathbf{B})\) 中在各自纤维内孤立的点集是 \(\operatorname{Spec}(\mathbf{B})\) 的开子集。
\end{enumerate}

\begin{enumerate}
\def\labelenumi{\arabic{enumi})}
\setcounter{enumi}{11}
\tightlist
\item
  设 H 为有限生成的分次代数，\(H_{0}\) 是一个域，\((\xi_{a})_{a\in\mathbf{A}}\) 是生成 H 的齐次元素。假设存在严格正整数族 \((d_{i})_{i\in\mathbf{I}}\)，使得 \(P_{H}=\prod(1-T^{d_{i}})^{-1}\)。以 \(\delta_{a}>0\) 表示 \(\xi_{a}\) 的次数。
\end{enumerate}

\begin{enumerate}
\def\labelenumi{\alph{enumi})}
\item
  证明存在单射 \(\varphi: I \to A\)，使得对于每个 \(i\)，\(d_i\) 整除 \(\delta_{\varphi(i)}\)。（可以使用 § 4, no. 2, 定理 1。）
\item
  证明有 \(\inf_{a \in A} \delta_a \leqslant \inf_{i \in I} d_i\)。
\item
  由此推出，若各 \(\delta_{a}\) 彼此相等，则分次代数 H 是正则的。
\item
  假设 A 有两个元素 \(a\) 和 \(a'\)，使得 \(\delta_{a} < \delta_{a'}\)，且 \(\delta_{a'}\) 不是 \(\delta_{a}\) 的倍数。证明分次代数 H 是正则的。
\end{enumerate}

\exercisesection{重数}

\begin{enumerate}
\def\labelenumi{\arabic{enumi})}
\tightlist
\item
  设 \(\Gamma \subset \mathbf{N}\) 为 \(\mathbf{N}\) 的子幺半群，使得 \(\mathbf{N} - \Gamma\) 是有限集。设 \((a_1, \ldots, a_r)\)（其中 \(0 < a_1 < \cdots < a_r\)）为 \(\Gamma\) 的极小生成元系，且 \(k\) 为一个域。置 \(\mathbf{A} = k[\Gamma]\)。
\end{enumerate}

\begin{enumerate}
\def\labelenumi{\alph{enumi})}
\item
  证明 A 同构于 \(k[\mathbf{T}]\) 中由 \(\mathbf{T}^{a_1}\)、\(\ldots\)、\(\mathbf{T}^{a_r}\) 生成的 k-子代数。
\item
  设 m 为由 \(T^{a_{1}}\)、\(\ldots\)、\(T^{a_{r}}\) 生成的 A 的极大理想；设 \(A_{m}\) 为 A 在 m 处的局部环，并置 \(q = mA_{m}\)。证明有 \(\dim_{k}q/q^{2} = r\) 及 \(e_{q}(A_{m}) = a_{1}\)。
\end{enumerate}

\begin{enumerate}
\def\labelenumi{\arabic{enumi})}
\setcounter{enumi}{1}
\tightlist
\item
  设 \(k\) 为一个域，\(m\) 为满足 \(\geqslant 1\) 的整数，I 为多项式环 \(k[\mathbf{T}_{1},\ldots,\mathbf{T}_{m}]\) 的一个理想，它由 \(\mathbf{M}_{1},\ldots,\mathbf{M}_{s}\) 所表示的单项式生成，这些点属于 \(\mathbf{R}_{+}^{m}\)，即 \(\mathbf{R}^{m}\) 的“正象限”。记 \(\mathcal{N}\) 为 \(\mathbf{R}_{+}^{m}\) 中的集合 \(\mathbf{M}_{i}+\mathbf{R}_{+}^{m}\)（\(1\leqslant i\leqslant s\)）之并的凸包。
\end{enumerate}

\begin{enumerate}
\def\labelenumi{\alph{enumi})}
\tightlist
\item
  置 \(m = (T_1, \ldots, T_m)\)、\(A = k[T_1, \ldots, T_m]_m\) 及 \(q = I.A\)。证明 \(A/q\) 具有有限长度，当且仅当 \(\mathbf{R}_+^m - \mathcal{N}\) 的体积有限。
\item
  证明等式 \(e_q(A) = m! \operatorname{Vol}(\mathbf{R}_+^m - \mathcal{N})\)。
\end{enumerate}

\begin{enumerate}
\def\labelenumi{\arabic{enumi})}
\setcounter{enumi}{2}
\tightlist
\item
  设 \(k\) 为一个域，\(f_{1},...,f_{r}\) 为分次环 \(k[X_{1},...,X_{n}]\) 中构成完全相交序列的齐次元素。称商环
\end{enumerate}

\[
\mathbf {H} = k [ \mathrm{X} _ {1},..., \mathrm{X} _ {n} ] / (f _ {1},..., f _ {r})
\]

  是分次完全相交环。设 I 为环 \(\mathbf{H} = k[\mathbf{X}_1, \ldots, \mathbf{X}_n] / (f_1, \ldots, f_r)\) 的非零分次理想，并赋予其商分次。证明有 \(\dim(\mathbf{I}) = \dim(\mathbf{H})\)。推出或者 \(\dim(\mathbf{H}/\mathbf{I}) < \dim(\mathbf{H})\)，或者 \(\dim(\mathbf{H}/\mathbf{I}) = \dim(\mathbf{H})\) 且 \(c_{\mathbf{M}/\mathbf{I}} < c_{\mathbf{M}}\)（\(\S 4, \text{n}^{\circ} 2\)）。（考虑由 \(s \in \mathbf{H}\) 中满足 \(s.\mathbf{I} = \{0\}\) 的元素构成的理想 b，并证明若 \(\dim(\mathbf{I}) < \dim(\mathbf{H})\)，则 b 含有一个非零因子元素；证明方法是考察 b 相对于 H 的极小素理想及 \(\mathbf{H}_{+} = \mathbf{H}_{\geqslant 1}\) 的位置。））

\begin{enumerate}
\def\labelenumi{\arabic{enumi})}
\setcounter{enumi}{3}
\tightlist
\item
  设 A 为诺特局部环，极大理想为 m。
\end{enumerate}

\begin{enumerate}
\def\labelenumi{\alph{enumi})}
\item
  假设 \(\mathrm{gr}_{\mathfrak{m}}(\mathrm{A})\) 是分次完全相交环（参见习题 3）。证明 A 的完备化 \(\widehat{A}\) 是完全相交的，即是正则局部环对一个由完全相交序列生成的理想的商。
\item
  设 \(x \in \mathfrak{m}\)，且 \(\xi\) 为 \(x\) 在 \(\mathfrak{m}^{\nu}/\mathfrak{m}^{\nu+1}\) 中的像，其中 \(\nu\) 满足 \(x \in \mathfrak{m}^{\nu}\) 且 \(x \notin \mathfrak{m}^{\nu+1}\)。证明 \(e_{\mathfrak{m}}(\mathbf{A}/x\mathbf{A}) = \nu e_{\mathfrak{m}}(\mathbf{A})\) 当且仅当 \(\xi\) 不是 \(\mathrm{gr}_{\mathfrak{m}}(\mathbf{A})\) 中的零因子，并且在此情形下有 \(\mathrm{gr}_{\mathfrak{m}}(\mathbf{A}/x\mathbf{A}) = \mathrm{gr}_{\mathfrak{m}}(\mathbf{A})/\xi\mathrm{gr}_{\mathfrak{m}}(\mathbf{A})\)，且 \(x\) 不是 A 中的零因子。（令 I 为 \(\mathrm{gr}_{\mathfrak{m}}(\mathbf{A})\) 的齐次理想，由 \(\eta \in \mathrm{gr}_{\mathfrak{m}}(\mathbf{A})\) 中满足 \(\xi\eta = 0\) 的元素构成。考虑 A 的理想 \(\mathrm{N}_n = \mathfrak{m}^{\nu+n}:x\)，以及从 \(\mathrm{I}_n\) 到 \(\mathrm{N}_{n+1}/\mathfrak{m}^{n+1}\) 的映射，它将 \(\eta \in \mathrm{I}_n\) 送到模 \(\mathfrak{m}^{n+1}\) 下的、元素 \(y \in A\) 提升 \(\eta\) 后的类。我们将证明它是单射，并由此推出不等式
\end{enumerate}

\[
\operatorname{long} _ {\mathbf {A} / \mathfrak {m}} (\mathbf {I} _ {n}) \leqslant \operatorname{long} _ {\mathbf {A} / \mathfrak {m}} (\mathbf {N} _ {n + 1} / \mathfrak {m} ^ {n + 1});\tag{i}
\]

此外，我们将证明等式

\[
(1 - \mathrm{T} ^ {\mathrm{v}}) \mathrm{H} _ {\mathrm{A}} ^ {(1)} = \mathrm{T} ^ {- \mathrm{v}} \mathrm{H} _ {\mathrm{A} / x \mathrm{A}} ^ {(1)} - \sum_ {n \geqslant \mathrm{v}} \operatorname{long} _ {\mathrm{A} / \mathfrak {m}} (\mathrm{N} _ {n} / \mathfrak {m} ^ {n}). \mathrm{T} ^ {n}.\tag{ii}
\]

最后注意到，\(\xi\) 是零因子当且仅当 I \(\neq\) \{0\}；利用上一习题，由 (i) 推出 \(\sum_{n\geqslant 0}\mathrm{long}_{\mathrm{A / m}}(\mathbf{N}_{n + 1} / \mathfrak{m}^{n + 1})\) \(\mathbf{T}^n = \frac{\mathbf{S}(\mathbf{T})}{(1 - \mathbf{T})^\nu}\)，其中 \(\mathbf{S}(1)\neq 0\)，并由 (ii) 得到不等式 \(e_{\mathfrak{m}}(\mathbf{A} / x\mathbf{A}) > e_{\mathfrak{m}}(\mathbf{A})\)。c) 设 \((x_{1},\dots,x_{s})\) 为极大理想 m 中元素组成的序列。假设 \(x_{i}\in \mathfrak{m}^{\nu_i}\) 且 \(x_{i}\notin \mathfrak{m}^{\nu_i + 1}\)，并令 \(\xi_{i}\) 为 \(x_{i}\) 在 \(\mathfrak{m}^{\nu_i} / \mathfrak{m}^{\nu_i + 1}\) 中的类。证明有 \(e_{\mathfrak{m}}(\mathbf{A} / x)\geqslant e_{\mathfrak{m}}(\mathbf{A})\nu_{1}\dots \nu_{s}\)，且等号成立当且仅当序列 \((\xi_1,\dots,\xi_s)\) 在环 \(\mathrm{gr}_{\mathfrak{m}}(\mathbf{A})\) 中是完全相交的。在此情形下，\(\mathrm{gr}_{\mathfrak{m}}(\mathbf{A} / x)\) 与 \((\mathrm{gr}_{\mathfrak{m}}(\mathbf{A}))/\xi\) 之间有分次环同构，其中 \(\mathbf{x} = \sum_{i}x_{i}\mathbf{A}\)，\(\xi = \sum_{i}\xi_{i}\mathrm{gr}_{\mathfrak{m}}(\mathbf{A})\)。

\begin{enumerate}
\def\labelenumi{\arabic{enumi})}
\setcounter{enumi}{4}
\tightlist
\item
  设 \(k\) 为一个域。置 \(\mathbf{R} = k[[\mathrm{X},\mathrm{Y},\mathrm{Z}]]\)，并考虑极大理想为 m 的局部环 \(\mathbf{A} = \mathbf{R}/(\mathbf{XY},\mathbf{X}\mathbf{Z})\)。证明有
\end{enumerate}

\[
\mathrm{H} _ {\mathrm{A}} = (1 - \mathrm{T}) (1 - \mathrm{T} - \mathrm{T} ^ {2}) \mathrm{H} _ {\mathrm{R}} = 1 + \sum_ {n \geqslant 1} (n + 2) \mathrm{T} ^ {n}.
\]

因此有 \(\dim(A) = 2\)、\(\dim_k(m/m^2) = 3\) 及 \(e_m(A) = 1\)。（我们有 \(H_A = P_G\)，其中 \(G = gr_m(A) = k[X, Y, Z]/(XY, XZ)\)。我们将证明存在分次 H-模的正合序列

\[
0 \rightarrow \mathrm{H} (- 3) \rightarrow \mathrm{H} (- 2) \oplus \mathrm{H} (- 2) \rightarrow \mathrm{H} \rightarrow \mathrm{G} \rightarrow 0
\]

\(\text{其中 } \mathrm{H} = k[\mathrm{X},\mathrm{Y},\mathrm{Z}].\)

\begin{enumerate}
\def\labelenumi{\arabic{enumi})}
\setcounter{enumi}{5}
\item
  设 \(k\) 为一个域，A 为维数 1 的局部、完备、约化 \(k\)-代数，\(m\) 为其极大理想。证明 A 的一个理想的最小生成元数至多为 \(e_m(A)\)。存在 \(X \in A\)，使得从 \(k[[X]]\) 到 A 的单射使 A 成为 \(k[[X]]\)-自由模，秩为 \(e_m(A)\)。
\item
  设 A 为极大理想为 m 的局部环，M 为有限生成 A-模，\((x_{1}, \ldots, x_{d})\) 为 m 中元素组成的、对 A-模 M 相交的序列，x 为理想 \(x_{1}A + \cdots + x_{d}A\)。
\end{enumerate}

\begin{enumerate}
\def\labelenumi{\alph{enumi})}
\tightlist
\item
  假设 A-模 M 是忠实的。则 \(\mathrm{gr}_{x}(\mathbf{A})\) 除以其幂零根基 n 所得的商同构于 \((\mathbf{A}/\mathfrak{m})[\mathbf{X}_{1}, \ldots, \mathbf{X}_{d}]\)，且 \((\mathrm{gr}_{x}(\mathbf{M}))_{n}\) 是一个 \((\mathrm{gr}_{x}(\mathbf{A}))_{n}\)-模，有限长度为 \(e_{x}(\mathbf{M})\)。
\item
  一般情形下，\(\mathrm{gr}_{x}(\mathbf{A})\) 有唯一的极小素理想 n，且 \((\mathrm{gr}_{x}(\mathbf{M}))_{n}\) 是一个 \(\mathrm{gr}_{x}(\mathbf{A})\)-模，有限长度为 \(e_{x}(\mathbf{M})\)。
\end{enumerate}

\begin{enumerate}
\def\labelenumi{\arabic{enumi})}
\setcounter{enumi}{7}
\tightlist
\item
  设 A 为一个环，a 为 A 的一个理想。置 \(a^{i} = A\) 对 \(i \leqslant 0\) 成立，并考虑 A-代数 \(\mathcal{A}(\mathfrak{a}) = \sum_{i \in \mathbf{Z}} \mathfrak{a}^{-i} v^{i} \subset \mathrm{A}[v, v^{-1}]\) 及 A-代数 \(\mathrm{P}(\mathfrak{a}) = \sum_{i \geqslant 0} \mathfrak{a}^{i} v^{i} \subset \mathrm{A}[v]\)。这两个分次 A-代数具有相同的全分式环。证明对于元素 \(a \in A\)，下列条件等价：
\end{enumerate}

\begin{enumerate}
\def\labelenumi{(\roman{enumi})}
\item
  元素 \(av\) 属于 \(\mathcal{A}(\mathfrak{a})\)-代数 \(\mathbf{A}[v, v^{-1}]\)，且它在 \(\mathcal{A}(\mathfrak{a})\) 上整。
\item
  P(a)-代数 A{[}v{]} 的元素 av 在 P(a) 上整。
\item
  存在整依赖关系：
\end{enumerate}

\[
a ^ {k} + b _ {1} a ^ {k - 1} + \dots + b _ {k} = 0 \quad \text { with } \quad b _ {j} \in \mathfrak {a} ^ {j}.
\]

称 \(a\) 关于理想 \(\mathfrak{a}\) 整。

\begin{enumerate}
\def\labelenumi{\arabic{enumi})}
\setcounter{enumi}{8}
\tightlist
\item
  \begin{enumerate}
  \def\labelenumii{\alph{enumii})}
  \tightlist
  \item
    设 A 和 a 如上，b 为 A 的一个包含 a 的理想。证明下列条件等价：
  \end{enumerate}
\end{enumerate}

\begin{enumerate}
\def\labelenumi{(\roman{enumi})}
\item
  b 的每个元素都关于 a 整。
\item
  \(\mathcal{A}(\mathfrak{a})\)-分次代数 \(\mathcal{A}(\mathfrak{b})\) 是整的。
\item
  分次 \(P(\mathfrak{a})\)-代数 \(P(\mathfrak{b})\) 是整的。
\item
  存在整数 \(k \geqslant 0\)，使得 \(\mathfrak{a}\mathfrak{b}^{k} = \mathfrak{b}^{k+1}\)，从而有 \(\mathfrak{a}^{n}\mathfrak{b}^{k} = \mathfrak{b}^{k+n}\) 对所有正整数 \(n\) 成立。称 \(\mathfrak{b}\) 关于 \(\mathfrak{a}\) 整。
\end{enumerate}

\begin{enumerate}
\def\labelenumi{\alph{enumi})}
\setcounter{enumi}{1}
\item
  证明在关于 a 整的理想中存在最大的理想。记其为 \(\overline{a}\)，称为 a 在 A 中的整闭包。验证若 \(a \subset b\)，则 \(\overline{a} \subset \overline{b}\)，且 \(\overline{\overline{a}} = \overline{a}\)。证明若整闭包 \(\overline{P(a)}\) 是有限生成的 \(P(a)\)-模，则当 k 充分大时，\(\mathfrak{a}\,\overline{\mathfrak{a}^{k}} = \overline{\mathfrak{a}^{k+1}}\)。
\item
  证明若 \(a \in A\) 满足整依赖关系
\end{enumerate}

\[
a ^ {k} + b _ {1} a ^ {k - 1} + \dots + b _ {k} = 0 \quad \text { with } \quad k \geqslant 1 \quad \text { and } \quad b _ {i} \in \mathfrak {a} ^ {i},
\]

 则 \(\mathfrak{a}(\mathfrak{a} + \mathbf{A}a)^{k-1} = (\mathfrak{a} + \mathbf{A}a)^k\)。由此推出，\(\mathfrak{a}\) 的整闭包在 A 中恰好是 \(\mathbf{A}\) 中关于 \(\mathfrak{a}\) 整的元素集合。

10) 设 A 为一个环。

\begin{enumerate}
\def\labelenumi{\alph{enumi})}
\item
  设 B 为整的 A-代数。证明对于 A 的每个元素 a 及每个理想 \(\mathfrak{a}\)，元素 a 关于 \(\mathfrak{a}\) 整，当且仅当 a 关于 B 中理想 \(\mathfrak{a}B\) 整，该理想由 \(\mathfrak{a}\) 生成。
\item
  假设 A 是诺特约化环，记 \(\mathfrak{p}_{1},\ldots,\mathfrak{p}_{r}\) 为 A 的极小素理想，并考虑自然单射：
\end{enumerate}

\[
\mathrm{A} \rightarrow \prod_ {1 \leqslant i \leqslant r} \mathrm{A} / \mathfrak {p} _ {i}.
\]

证明它使 \(\prod_{1\leqslant i\leqslant r}\mathrm{A} / \mathfrak{p}_i\) 成为整的 A-代数。

\begin{enumerate}
\def\labelenumi{\alph{enumi})}
\setcounter{enumi}{2}
\item
  由此推出，给定诺特环 A，A 的元素 a 关于理想 \(\mathfrak{a}\) 整的充要条件是：对于 A 的每个极小素理想 p，a 在 A/p 中的像关于理想 \((\mathfrak{a} + p)/p\) 整。
\item
  证明若 \(\overline{a} = \overline{b}\)，则理想 \(a\) 和 \(b\) 具有相同的极小相伴素理想。
\end{enumerate}

\begin{enumerate}
\def\labelenumi{\arabic{enumi})}
\setcounter{enumi}{10}
\tightlist
\item
  设 A 为诺特环，a、b 为 A 的两个理想，且 Supp(a) = Supp(b) = Spec(A)。证明下列条件等价：
\end{enumerate}

\begin{enumerate}
\def\labelenumi{(\roman{enumi})}
\tightlist
\item
  存在有限生成的 A-模 M，使得 \(\operatorname{Supp}(\mathbf{M}) = \operatorname{Spec}(\mathbf{A})\)，且 \(bM \subset aM\)。\\
\item
  有 \(a, b \subset \overline{a}\)。
\end{enumerate}

\begin{enumerate}
\def\labelenumi{\arabic{enumi})}
\setcounter{enumi}{11}
\tightlist
\item
  设 A 为诺特环，a、b、c 为 A 的三个理想，且 c 包含于 A 的根基中。
\end{enumerate}

\begin{enumerate}
\def\labelenumi{\alph{enumi})}
\item
  证明若有包含关系 \(\overline{b + c.a} \subset \overline{a}\)，则有 \(\overline{b} \subset \overline{a}\)。
\item
  因此，若 A 是局部环，且有 \(b \subset a\) 和 \(\overline{b} = \overline{a}\)，则至少存在一个理想 \(b_{1} \subset b\)，使得
\end{enumerate}

\[
\overline {{{\mathbf {b} _ {1}}}} = \overline {{{\mathbf {a}}}},
\]

\(\beta)\) \(\mathbf{b}_{1}\) 对此性质是极小的，即若 \(\mathbf{b}_{1}^{\prime}\subset \mathbf{b}_{1}\) 且 \(\overline{\mathbf{b}_1'} = \overline{\mathbf{b}_1}\)，则 \(\mathbf{b}_1' = \mathbf{b}_1\)。

\begin{enumerate}
\def\labelenumi{\arabic{enumi})}
\setcounter{enumi}{12}
\item
  设 A 为局部环，a、b 为 A 中异于 A 的两个理想。证明若 a ⊂ b，则 \(\overline{a} = \overline{b}\) 当且仅当分次代数 gr \(_b(A)\) 在分次代数 gr \(_a(A)\) 上整。推出若 a = (x \(_1\) , \ldots, x \(_d\) ) ⊂ b，其中 d = dim(A)，则 \(\overline{a} = \overline{b}\) 当且仅当初始形式 \(\xi_{1}\) , \ldots, \(\xi_{d}\)（它们是在分次代数 gr \(_b(A)\) 中的 x \(_i\) 的初始形式）次数为 1，且构成 gr \(_b(A)\) 的极大相交序列。
\item
  设 A 为局部环，m 为其极大理想，q 为 A 中异于 A 的理想，且 A/q 具有有限长度。假设 A 含有子域 k，使得 A = k + m。
\end{enumerate}

\begin{enumerate}
\def\labelenumi{\alph{enumi})}
\tightlist
\item
  证明 q 中存在元素 \(x_{1}, \ldots, x_{d}\)，其中 \(d = \dim(A)\)，且 \(x_{i} \in q^{\delta_{i}}\)，使得对于每个充分大的整数 v，有 \(q^{v} = x_{1} \cdot q^{v - \delta_{1}} + \cdots + x_{d} \cdot q^{v - \delta_{d}}\)。
\item
  若假设域 A/m 是无限的，则可以找到一个由 q 中 d 个元素生成的理想 a，其中 \(d = \dim(A)\)，使得 \(\overline{a} = \overline{q}\)（可以使用 § 7, no. 5。）
\end{enumerate}

\begin{enumerate}
\def\labelenumi{\arabic{enumi})}
\setcounter{enumi}{14}
\tightlist
\item
  设 A 为极大理想 m 的局部、诺特、完备环。置 \(d = \dim(A)\)。假设 A 含有子域 k，使得 \(A = k + m\)。
\end{enumerate}

证明可以在 m 中找到 \(x_{1},\ldots,x_{d}\)，使得：

  态射 \(\varphi:k[[X_{1},...,X_{d}]]\to A\) 由 \(\varphi(X_{i})=x_{i}\) 定义，使 A 成为整的 \(k[[X_{1},...,X_{d}]]\)-代数；

\(\beta)\) 记 \(\mathfrak{m}_0 = (X_1, \ldots, X_d)\) 为 \(k[[X_1, \ldots, X_d]]\) 的极大理想，则有 \(\overline{\mathfrak{m}_0A} = \mathfrak{m}\)，因而 \(e_{\mathfrak{m}_0\mathbf{A}}(\mathbf{A}) = e_{\mathfrak{m}}(\mathbf{A})\)。

  证明若 \(x_{1},\ldots,x_{d}\) 满足条件 \(\alpha\)，则条件 \(\beta\) 成立当且仅当初始形式 \(\xi_{1},\ldots,\xi_{d}\)（它们是 \(x_{i}\) 在 \(\mathrm{gr}_{\mathfrak{m}}(\mathbf{A})\) 中的初始形式）构成极大相交序列。

\begin{enumerate}
\def\labelenumi{\arabic{enumi})}
\setcounter{enumi}{15}
\item
  设 A 为诺特局部环，\(x = (x_1, \ldots, x_d)\) 为由 A 的极大相交序列生成的理想。证明有不等式 \(e_x(A) \leqslant \text{long}(A/x)\)，且等号成立只能在同态 \((A/x)[T_1, \ldots, T_d] \xrightarrow{\varphi} \text{gr}_x(A)\) 为同构时成立；该同态由 \(\varphi(T_i) = \xi_i\) 定义，其中 \(\xi_i\) 是 \(x_i\) 模 \(x^2\) 的类（这意味着 \((x_1, \ldots, x_d)\) 是完全相交的）。
\item
  设 A 为诺特局部环，具有长度 \(d = \dim(A)\) 的完全相交序列；令 \(\mathfrak{a}\) 为由此序列生成的理想，令 \(f \in A\) 为关于 \(\mathfrak{a}\) 整但不属于该理想的元素（例如 \(A = k[[X_1, ..., X_d]]\)，其中 k 为域，\(\mathfrak{a} = (X_1^k, ..., X_d^k)\)，且 \(f = X_1^{k-1}X_2\)）。令 \(\mathfrak{a}' = \mathfrak{a} + fA\)。证明有不等式 \(e_{\mathfrak{a}'}(A) > \text{long}(A/\mathfrak{a}')\)。
\item
  设 A 为一个环，\(\mathfrak{a}\) 为 A 的一个理想。对于每个元素 \(a \in \mathbf{A}\)，以 \(\nu_{\mathfrak{a}}(a)\) 表示整数 \(\nu \in \mathbf{N}\) 中满足 \(a \in \mathfrak{a}^{\nu}\) 者的集合的上界。
\end{enumerate}

\begin{enumerate}
\def\labelenumi{\alph{enumi})}
\tightlist
\item
  证明若 \(a \notin \bigcap_{k=0}^{\infty} \mathfrak{a}^k\)，则序列 \(\left(\frac{\nu_{\mathfrak{a}}(a^k)}{k}\right)\) 收敛。以 \(\overline{\nu}_{\mathfrak{a}}(a)\) 表示其极限。若
\end{enumerate}

\(a \in \bigcap_{k=0}^{\infty} \mathfrak{a}^{k}\)，则置 \(\overline{\nu}_{\mathfrak{a}}(a) = +\infty\)。

\begin{enumerate}
\def\labelenumi{\alph{enumi})}
\setcounter{enumi}{1}
\tightlist
\item
  证明对于所有 \(a \in \mathbf{A}\)，有 \(\overline{\nu}_{\mathfrak{a}}(a) = \overline{\nu}_{\mathfrak{a}}(a)\)，并由此推出下列条件等价：
\end{enumerate}

\begin{enumerate}
\def\labelenumi{(\roman{enumi})}
\item
  \(\overline{\nu}_{\mathfrak{a}}(a)\geqslant 1,\)
\item
  \(a \in \overline{\mathfrak{a}}\).
\end{enumerate}

（研究 a) 中所考察序列的单调性。）

\begin{enumerate}
\def\labelenumi{\arabic{enumi})}
\setcounter{enumi}{18}
\tightlist
\item
  设 A 为诺特环，q 为 A 的一个理想，且 A/q 具有有限长度，q 包含于 A 的根基中。证明有
\end{enumerate}

\[
e _ {\mathfrak {q}} (\mathrm{A}) = e _ {\overline {{\mathfrak {q}}}} (\mathrm{A}).
\]

\begin{enumerate}
\def\labelenumi{\arabic{enumi})}
\setcounter{enumi}{19}
\tightlist
\item
  设 A 为诺特环，M 为有限生成的 A-模，\(q_{1}, \ldots, q_{k}\) 为 A 的理想，且对于每个 i，\(M/q_{i}M\) 具有有限长度。假设 \(M \neq 0\)，并置 \(d = \dim_{\mathbf{A}}(\mathbf{M})\)。a) 证明 \(M/q_{1}^{n_{1}} \ldots q_{k}^{n_{k}}M\) 对每个由正整数组成的 k-元组 \((n_{1}, \ldots, n_{k})\) 都具有有限长度。
\end{enumerate}

\begin{enumerate}
\def\labelenumi{\alph{enumi})}
\setcounter{enumi}{1}
\tightlist
\item
  考虑类型为 \(\mathbf{N}^k\) 的分次环 H，使得
\end{enumerate}

\[
\mathrm{H} _ {n _ {1}, \dots , n _ {k}} = \mathfrak {q} _ {1} ^ {n _ {1}} \mathfrak {q} _ {2} ^ {n _ {2}} \dots \mathfrak {q} _ {k} ^ {n _ {k}} / \mathfrak {q} _ {1} ^ {n _ {1} + 1} \mathfrak {q} _ {2} ^ {n _ {2}} \dots \mathfrak {q} _ {k} ^ {n _ {k}},
\]

以及分次 H-模 N，使得

\[
 \mathrm{N} _ {n _ {1}, \dots , n _ {k}} = \mathrm{q} _ {1} ^ {n _ {1}} \mathrm{q} _ {2} ^ {n _ {2}} \dots \mathrm{q} _ {k} ^ {n _ {k}} \mathrm{M} / \mathrm{q} _ {1} ^ {n _ {1} + 1} \mathrm{q} _ {2} ^ {n _ {2}} \dots \mathrm{q} _ {k} ^ {n _ {k}} \mathrm{M};
\]

证明 N 是有限生成的分次 H-模，由 \(N_{0,\ldots,0}\) 及有限个次数为 \(\mathbf{Z}^{k}\) 的标准基向量的元素生成。

由此推出，存在多项式 \(\mathbf{Q}(\mathrm{T}) \in \mathbf{Z}[\mathrm{T}_1,..,\mathrm{T}_k]\) 及整数 \(s_1 \geqslant 1, ..., s_k \geqslant 1\)，使得有

\[
\sum_ {n \in \mathbb {N} ^ {k}} \operatorname{long} _ {\mathrm{A}} (\mathrm{H} _ {n}). \mathrm{T} ^ {n} = \frac {\mathrm{Q} (\mathrm{T})}{(1 - \mathrm{T} _ {1}) ^ {s _ {1}} \dots (1 - \mathrm{T} _ {k}) ^ {s _ {k}}}.
\]

（我们对 d 使用归纳法。）

\begin{enumerate}
\def\labelenumi{\alph{enumi})}
\setcounter{enumi}{2}
\tightlist
\item
  由此，通过将 § 6, no. 1 的结果延拓到 \(\mathbf{Z}((\mathrm{T}_{1},\ldots,\mathrm{T}_{k}))\)，推出存在整数 \(c(\alpha_{1},\ldots,\alpha_{k})\)，对应于序列 \((\alpha_{1},\ldots,\alpha_{k})\)，其中 \(|\alpha| = \sum_{i=1}^{k}\alpha_{i} = d\)，使得有
\end{enumerate}

\[
e _ {\mathfrak {q} _ {1} ^ {\nu_ {1}} \dots \mathfrak {q} _ {k} ^ {\nu_ {k}}} (\mathbf {M}) = \sum_ {| \alpha | = d} \frac {d !}{\alpha_ {1} ! \dots \alpha_ {k} !} c (\alpha_ {1},..., \alpha_ {k}) v _ {1} ^ {\alpha_ {1}}... v _ {k} ^ {\alpha_ {k}}
\]

其中 \((\mathbf{v}_1,\dots,\mathbf{v}_k)\) 属于 \(\mathbf{N}^k\)

\begin{enumerate}
\def\labelenumi{\alph{enumi})}
\setcounter{enumi}{3}
\tightlist
\item
  现在假设 A 是剩余域为无限域 κ 的局部环。利用分次 H-模 N 的表面元素的存在性，证明有
\end{enumerate}

\[
c (\alpha_ {1}, \dots , \alpha_ {k}) = \inf _ {\mathbf {q} _ {\alpha}} (e _ {\mathbf {q} _ {\alpha}} (\mathbf {M}))
\]

  其中 \(q_{\alpha}\) 遍历 A 的理想集，这些理想分别由 \(\alpha_{1}\) 个 \(q_{1}\) 中的元素、\(\alpha_{2}\) 个 \(q_{2}\) 中的元素、\ldots、\(\alpha_{k}\) 个 \(q_{k}\) 中的元素生成。有时我们将 \(e(\mathbf{q}_{1}^{[\alpha_{1}]} + \cdots + \mathbf{q}_{k}^{[\alpha_{k}]}; \mathbf{M})\) 记为整数 \(c(\alpha_{1}, \ldots, \alpha_{k})\)；当 M = A 时，简写为 \(e(\mathbf{q}_{1}^{[\alpha_{1}]} + \cdots + \mathbf{q}_{k}^{[\alpha_{k}]})\)。e) 证明等式

\[
e _ {\mathfrak {q} _ {1}, \mathfrak {q} _ {2}} (\mathbf {M}) = \sum_ {i = 1} ^ {d} \binom {d} {i} e (\mathfrak {q} _ {1} ^ {[ i ]} + \mathfrak {q} _ {2} ^ {[ d - i ]}; \mathbf {M}).
\]

\begin{enumerate}
\def\labelenumi{\alph{enumi})}
\setcounter{enumi}{5}
\tightlist
\item
  证明有：
\end{enumerate}

\[
e \left(\overline {{{\mathfrak {q}}}} _ {1} ^ {[ \alpha_ {1} ]} + \dots + \overline {{{\mathfrak {q}}}} _ {k} ^ {[ \alpha_ {k} ]}; \mathrm{M}\right) = e \left(\mathfrak {q} _ {1} ^ {[ \alpha_ {1} ]} + \dots + \mathfrak {q} _ {k} ^ {[ \alpha_ {k} ]}; \mathrm{M}\right)
\]

其中 \(\overline{\mathfrak{q}}_{i}\) 表示理想 \(\mathfrak{q}_{i}\) 在 A 中的整闭包（参见习题 19）。

\begin{enumerate}
\def\labelenumi{\arabic{enumi})}
\setcounter{enumi}{20}
\tightlist
\item
  设 A 和 A' 为两个诺特环，q（分别为 q'）为 A（分别为 A'）的理想，包含于 A（分别为 A'）的根基中，且 \(\operatorname{long}(A/q)\)（分别为 \(\operatorname{long}(A'/q'))\) 有限。设 Q 为理想 \(q \otimes A' + A \otimes q'\)，它属于环 \(A \otimes_{\mathbf{Z}} A'\)。证明它包含于 \(A \otimes_{\mathbf{Z}} A'\) 的根基中，并且有等式
\end{enumerate}

\[
e _ {\mathfrak {Q}} (\mathrm{A} \otimes_ {\mathbf {Z}} \mathrm{A} ^ {\prime}) = e _ {\mathfrak {q}} (\mathrm{A}). e _ {\mathfrak {q} ^ {\prime}} (\mathrm{A} ^ {\prime}).
\]

22) * 设 A 为诺特整环，且为局部完备环（分别为域 k 上的有限生成代数），\(a \in A\) 为非可逆非零元素。作下列假设：

α) \(\operatorname{Supp}(A/aA)\) 有唯一极小元 p。

\(\beta)\) 有 \(aA_{p} = pA_{p}\)。

γ) A/p 是整闭环。

本习题旨在证明，对于 A 中每个包含 a 的极大理想 m，环 \(A_{m}\) 是整闭的，且有 p = aA（“Hironaka 引理”）。本习题将使用串联性。在第 X 章中将看到，整的诺特局部完备环满足此性质。域上的有限型整代数也满足此性质（§ 2, n° 4, 定理 3）。

\begin{enumerate}
\def\labelenumi{\alph{enumi})}
\item
  证明 \(A_p\) 是离散赋值环（§ 3, n° 1, 命题 1 及 VI, § 3, n° 6, 命题 9）。
\item
  作 \(\alpha\) 和 \(\beta\) 的假设，并假设 A 是整闭的。证明有 \(\mathfrak{p} = a\mathbf{A}\)（VII, § 1, n° 4, 命题 8）。
\item
  设 A' 为 A 的整闭包。证明 A' 是整的诺特局部完备环（分别证明它是 k 上的有限型代数）。（注意 A 是日式环（IX, § 4, n° 2, 定理 2），因而 A' 是整的、诺特的、半局部且完备的，从而是局部的（III, § 2, n° 13, 命题 19 的推论）。当 A 是 k 上的有限型代数时，使用 V, § 3, n° 2, 定理 2。）
\item
  证明有 A' ⊗A A\_p = A\_p（使用 a) 及 V, § 1, n° 5, 命题 16）。推出存在 A' 中位于 p 之上的唯一素理想 p'，且有 A\_p = A'\_p'、aA'\_p' = p'A'\_p'。
\item
  设 q 为 A'/aA' 的支集中的极小元。证明 q 是 A' 中位于 p 之上的素理想。（注意 q 的高度为 1（§ 3, n° 1, 命题 1），因而根据串联性有 \(\dim(A'/q)=\dim(A)-1\)。推出 \(\dim(A/(q \cap A))=\dim(A)-1\)，从而 \(q \cap A = p\)。）
\item
  从 e) 推出 \(aA'\) 是 \(A'\) 的素理想，且有 \(aA' = p' = pA'\)（可以使用 b)）。
\end{enumerate}

g) 证明有 \(\mathbf{A}_{\mathfrak{m}}^{\prime} = \mathbf{A}_{\mathfrak{m}}\) 对 A 中每个包含 \(a\) 的极大理想 m 成立。（由包含关系 \(\mathbf{A} \subset \mathbf{A}' \subset \mathbf{A}_{\mathfrak{p}}\) 推出包含关系 \(\mathbf{A}/\mathfrak{p} \subset \mathbf{A}'/\mathfrak{p}' \subset \mathbf{A}_{\mathfrak{p}}/\mathfrak{p}\mathbf{A}_{\mathfrak{p}}\)。推出 \(\mathbf{A}/\mathfrak{p}\) 和 \(\mathbf{A}'/\mathfrak{p}'\) 具有相同的分式域，从而根据 \(\gamma)\)，有 \(\mathbf{A}/\mathfrak{p} = \mathbf{A}'/\mathfrak{p}' = \mathbf{A}' \otimes_{\mathbf{A}} (\mathbf{A}/\mathfrak{p})\)，最后一个等式由 \(f)\) 得出。因此有 \((\mathbf{A}'/\mathbf{A}) \otimes_{\mathbf{A}} (\mathbf{A}/\mathfrak{p}) = 0\)。得出结论。）然后完成 Hironaka 引理的证明。 *

\begin{enumerate}
\def\labelenumi{\arabic{enumi})}
\setcounter{enumi}{22}
\tightlist
\item
  \begin{enumerate}
  \def\labelenumii{\alph{enumii})}
  \tightlist
  \item
    设 A 为满足 \(e(A) = 1\) 的诺特局部环。证明 A 有唯一的极小素理想 \(p\)，使得 \(\dim(A) = \dim(A/p)\)。此外证明有 \(\text{long}(A_p) = 1\) 及 \(e(A/p) = 1\)（\S 7, no. 1, 注记 4）。
  \end{enumerate}
\end{enumerate}

\begin{enumerate}
\def\labelenumi{\alph{enumi})}
\setcounter{enumi}{1}
\tightlist
\item
  回忆（p. 82, 习题 10），若对于 A 的每个极小素理想 q 都有 \(\dim(A/q) = \dim(A)\)，则称局部环是等维的；若 A-模 A 没有嵌入相伴素理想，则称其没有嵌入素理想（IV, § 2, no. 3, 注记）。* 将在第 X 章看到，Macaulay 局部环的整商的完备化具有这两个性质。* 设 A 为无嵌入素理想的等维诺特局部环，且 \(e(A) = 1\)。证明 A 是整环。
\end{enumerate}

24) * 设 A 为诺特局部环，\(\widehat{A}\) 为其完备化。本习题旨在证明下列性质等价：

\begin{enumerate}
\def\labelenumi{(\roman{enumi})}
\item
  A 是正则环；
\item
  \(\hat{\mathbf{A}}\) 是等维的且没有嵌入素理想，并且 \(e(\mathbf{A}) = 1\)；
\item
  \(\hat{\mathbf{A}}\) 是整环且 \(e(\mathbf{A}) = 1\)。
\end{enumerate}

证明蕴含 (i) \(\Rightarrow\) (ii) \(\Rightarrow\) (iii)。注意 (iii) 蕴含等式 \(e(\hat{\mathbf{A}}) = 1\)，且为证明 (iii) \(\Rightarrow\) (i)，只需证明 \(\hat{\mathbf{A}}\) 是正则环。因此，为证明 (iii) \(\Rightarrow\) (i)，可以假设 A 完备，以下均如此。我们先处理剩余域 \(\kappa_{\mathrm{A}}\) 含有无穷多个元素的情形。对 A 的维数作归纳。考察 \(\dim(A) = 0\) 的情形，以下设 \(\dim(A) > 0\)。此时存在表面元素 \(x \in \mathrm{A}\)（\S 7, no. 5, 注记 4）。根据串联性（习题 22）及 \S 3, no. 1, 命题 1，对于包含 xA 的每个极小素理想 p，都有 \(\dim(A/p) = \dim(A/xA)\)。由于有 \(e(\mathrm{A}/x\mathrm{A}) = 1\)（\S 7, no. 5, 定理 1），包含 xA 的极小素理想中存在唯一一个 p，且有 \(xA_{p} = pA_{p}\)、\(e(\mathrm{A}/p) = 1\)（习题 23）。根据归纳假设，A/p 是正则环，因而是整闭的（\S 5, no. 2, 定理 1 的推论 1）。由 Hironaka 引理（习题 22）推出 p = xA。因此 A/xA 是正则环，并且由于 A 是整环，A 是正则环（\S 5, no. 3, 命题 2）。

现在假设 \(\kappa_{A}\) 是有限域。证明只需证明爆破 A{]}X{[}（IX, 附录, no. 2）是正则环。注意 A 以及由此得到的 A{]}X{[} 同构于正则环的商（IX, § 2, no. 5, 定理 3），且 \(e(A]X[)=1\)。根据 Macaulay 环理论（第 X 章），完备化 \(\widehat{A]X[}\) 是等维的且没有嵌入素理想。推出 \(\widehat{A]X[}\) 是整环（习题 23）且 \(e(\widehat{A]X[})=1\)，因而 \(A]X[\) 是正则环。得出结论。*

\begin{enumerate}
\def\labelenumi{\arabic{enumi})}
\setcounter{enumi}{24}
\tightlist
\item
  设 \(k\) 为一个域，A 为 \(k\)-局部代数，即在有限型 \(k\)-代数的一个素理想处局部化所得。证明 A 是正则环当且仅当 A 是整环且 \(e(A) = 1\)。（可以参考习题 24 所用的方法，或者注意到由于 A 是正则环的整商，局部环 \(\hat{A}\) 是等维的且没有嵌入素理想。）
\end{enumerate}

\chapter{完备诺特局部环}

本章中所有环均假设为交换环；代数均为结合的、交换的且有单位元的。我们以 1 \(_{A}\) 表示环 A 的单位元。

若 A 为环且 p 为 A 的素理想，则以 \(\kappa(p)\) 表示局部环 \(A_p\) 的剩余域。若环 A 是局部环，则以 \(\mathfrak{m}_A\) 表示其极大理想，以 \(\kappa_A\) 或 \(\kappa(\mathfrak{m}_A)\) 表示其剩余域。

  若环同态 \(\rho: A \to B\) 使 B 成为平坦的（分别为忠实平坦的）A-模，则称该同态为平坦的（分别为忠实平坦的）。回忆（I, § 3, n° 5, 命题 9），若 A 和 B 是局部环，则 \(\rho\) 忠实平坦当且仅当它平坦且为局部同态。

\section{WITT 向量}

本节中始终以 p 表示素数。

\subsection{Witt 多项式}

  对于每个整数 \(n \geqslant 0\)，称元素 \(\Phi_n\) 为 \(\mathbf{Z}[X_0, ..., X_n]\) 中的第 \(n\) 个 Witt 多项式，其定义为

\[
\Phi_ {n} \left(\mathrm{X} _ {0}, \dots , \mathrm{X} _ {n}\right) = \sum_ {i = 0} ^ {n} p ^ {i} \mathrm{X} _ {i} ^ {p ^ {n - i}} = \mathrm{X} _ {0} ^ {p ^ {n}} + p \mathrm{X} _ {1} ^ {p ^ {n - 1}} + \dots + p ^ {n} \mathrm{X} _ {n}.\tag{1}
\]

显然有 \(\Phi_{0}=X_{0}\) 及递推关系

\begin{enumerate}
\def\labelenumi{(\arabic{enumi})}
\setcounter{enumi}{1}
\tightlist
\item
\end{enumerate}

\[
\Phi_ {n + 1} \left(\mathrm{X} _ {0}, \dots , \mathrm{X} _ {n + 1}\right) = \Phi_ {n} \left(\mathrm{X} _ {0} ^ {p}, \dots , \mathrm{X} _ {n} ^ {p}\right) + p ^ {n + 1} \mathrm{X} _ {n + 1}\tag{3}
\]

\[
\Phi_ {n + 1} \left(\mathrm{X} _ {0}, \dots , \mathrm{X} _ {n + 1}\right) = \mathrm{X} _ {0} ^ {p ^ {n + 1}} + p \Phi_ {n} \left(\mathrm{X} _ {1}, \dots , \mathrm{X} _ {n + 1}\right).
\]

将 \(X_{i}\) 的权赋为 \(p^{i}\) 时，多项式 \(\Phi_{n}\) 是权为 \(p^{n}\) 的等重多项式（A, IV, p.3）。

命题 1. --- 设 A 为滤过环，\((\mathbf{J}_n)_{n\in \mathbb{Z}}\) 为其滤过。假设 \(\mathbf{J}_0 = \mathbf{A}\) 且 \(p\cdot 1_{\mathbf{A}}\in \mathbf{J}_{1}\)。设 m、n 为满足 \(m\geqslant 1\) 及 \(n\geqslant 0\) 的整数，并设 \(a_0,\dots,a_n,b_0,\dots,b_n\) 为 A 的元素。

\begin{enumerate}
\def\labelenumi{\alph{enumi})}
\tightlist
\item
  若有 \(a_i \equiv b_i \mod J_m\) 对 \(0 \leqslant i \leqslant n\) 成立，则有
\end{enumerate}

\[
\Phi_ {i} (a _ {0}, \dots , a _ {i}) \equiv \Phi_ {i} (b _ {0}, \dots , b _ {i}) \bmod \mathrm{J} _ {m + i} \quad \text{当} \quad 0 \leqslant i \leqslant n.
\]

\begin{enumerate}
\def\labelenumi{\alph{enumi})}
\setcounter{enumi}{1}
\tightlist
\item
  假设对于每个整数 \(k \geqslant 1\) 及所有 \(x \in \mathbf{A}\)，关系 \(p x \in \mathbf{J}_{k+1}\) 蕴含 \(x \in \mathbf{J}_k\)。若有 \(\Phi_i(a_0, ..., a_i) \equiv \Phi_i(b_0, ..., b_i) \mod \mathbf{J}_{m+i}\) 对 \(0 \leqslant i \leqslant n\) 成立，则有 \(a_i \equiv b_i \mod \mathbf{J}_m\) 对 \(0 \leqslant i \leqslant n\) 成立。
\end{enumerate}

引理 1. --- 若 x、y 为 A 中模 \(J_{m}\) 同余的两个元素，则有

\[
x ^ {p ^ {n}} \equiv y ^ {p ^ {n}} \bmod J _ {m + n}.
\]

对 \(n\) 作归纳，将问题归约到 \(n=1\) 的情形。令 P 表示多项式 \(\sum_{i=0}^{p-1} X^i Y^{p-1-i}\)，它属于 Z{[}X, Y{]}。根据对 \(x\) 和 \(y\) 所作的假设，有 P(x, y) \(\equiv\) P(x, x) \(\equiv\) p, \(x^{p-1}\) mod. J \(_m\)。但 J \(_m\) + p.A \(\subset\) J \(_1\)，所以 P(x, y) \(\in\) J \(_1\)。最后，\(x^p - y^p = (x - y)\) P(x, y) 属于 J \(_m\) J \(_1\) \(\subset\) J \(_{m+1}\)。

我们对 n 作归纳来证明 a)。n = 0 的情形是直接的。假设 \(n \geqslant 1\)。在 a) 的假设下，有

\begin{enumerate}
\def\labelenumi{(\arabic{enumi})}
\setcounter{enumi}{3}
\tightlist
\item
\end{enumerate}

\(a_{i}^{p}\equiv b_{i}^{p}\bmod \mathrm{J}_{m + 1}\quad \text{当}\quad 0\leqslant i\leqslant n - 1\quad \text{由引理 } 1.\)

\[
\Phi_ {n - 1} (a _ {0} ^ {p}, \dots , a _ {n - 1} ^ {p}) \equiv \Phi_ {n - 1} (b _ {0} ^ {p}, \dots , b _ {n - 1} ^ {p}) \pmod{J _ {m + n}}\tag{5}
\]

根据归纳假设（应用于 A 中的元素 \(a_0^p, \ldots, a_{n-1}^p, b_0^p, \ldots, b_{n-1}^p\)），以及

\[
\Phi_ {n} (a _ {0},..., a _ {n}) - p ^ {n}. a _ {n} \equiv \Phi_ {n} (b _ {0},..., b _ {n}) - p ^ {n}. b _ {n} \pmod{J _ {m + n}}\tag{6}
\]

根据公式 (2) 和 (5)。由于 \(a_{n}-b_{n}\) 属于 \(\mathbf{J}_{m}\)，元素 \(p^{n}a_{n}-p^{n}b_{n}\) 属于 \(\mathbf{J}_{m+n}\)，由 (6) 可推出同余关系

\[
\Phi_ {n} (a _ {0}, \dots , a _ {n}) \equiv \Phi_ {n} (b _ {0}, \dots , b _ {n}) \pmod{J _ {m + n}},
\]

从而得到 a)。

  我们证明 \(b)\)，通过对 \(n\) 作归纳。\(n=0\) 的情形是直接的。假设 \(n\geqslant1\)。在 \(b)\) 的假设下，根据归纳假设，有 \(a_{i}\equiv b_{i}\) mod. \(\mathbf{J}_{m}\)（\(0\leqslant i\leqslant n-1\)），并像以前一样推出同余关系 (4)、(5) 和 (6)。但根据假设，\(\Phi_{n}(a_{0},...,a_{n})\) 与 \(\Phi_{n}(b_{0},...,b_{n})\) mod. \(\mathbf{J}_{m+n}\) 同余，因而有 \(p^{n}(a_{n}-b_{n})\in\mathbf{J}_{m+n}\)。由于关系 \(p x\in\mathbf{J}_{k+1}\) 蕴含 \(x\in\mathbf{J}_{k}\) 对每个 \(x\in\mathbf{A}\) 及每个 \(k\geqslant1\) 成立，故 \(a_{n}-b_{n}\in\mathbf{J}_{m}\)，证明完成。

\subsection{\texorpdfstring{映射 \(f, v\) 和 \(\Phi\)}{映射 f、v 和 Φ}}

  设 A 为一个环。在 \(\mathbf{A}^{\mathbf{N}}\) 上赋予乘积环结构。以 \(f_{\mathbf{A}}\) 或简称 \(f\) 表示自同态 \((a_{n})_{n\in\mathbf{N}}\mapsto(a_{n+1})_{n\in\mathbf{N}}\) 的 \(\mathbf{A}^{\mathbf{N}}\)。以 \(v_{\mathbf{A}}\) 或简称 \(v\) 表示 \(\mathbf{A}^{\mathbf{N}}\) 的底层加法群的自同态，它将 \((a_{n})_{n\in\mathbb{N}}\) 送到 \((0,p\cdot a_{0},p\cdot a_{1},...)\)。

  对于每个整数 \(m \geqslant 0\)，以 \(\Phi_{m}\) 表示从 \(\mathbf{A}^{\mathbf{N}}\) 到 A 的映射，它将 \(\boldsymbol{a} = (a_{n})_{n \in \mathbf{N}}\) 送到 \(\Phi_{m}(a_{0}, \ldots, a_{m})\)。以 \(\Phi_{\mathrm{A}}\) 或简称 \(\Phi\) 表示映射 \(\boldsymbol{a} \mapsto (\Phi_{n}(\boldsymbol{a}))_{n \in \mathbf{N}}\)，它从 \(\mathbf{A}^{\mathbf{N}}\) 到自身。

引理 2.--- 设 A 为带有自同态 σ 的环，对于每个 a ∈ A，σ(a) ≡ a\^{}p 模 p.A。设 n ≥ 1 为整数，a\_0, \ldots, a\_\{n-1\} 为 A 的元素。置 u\_i = Φ\_i(a\_0, \ldots, a\_i)，其中 0 ≤ i ≤ n - 1。设 u\_n 为 A 的一个元素。下列条件等价：

\begin{enumerate}
\def\labelenumi{\alph{enumi})}
\item
  存在 \(a_{n} \in A\)，使得 \(u_{n} = \Phi_{n}(a_{0}, \ldots, a_{n})\)。
\item
  \(\sigma(u_{n-1}) \equiv u_n \pmod{p^nA}\).
\end{enumerate}

对于 \(0 \leqslant i \leqslant n - 1\)，有 \(\sigma(a_i) \equiv a_i^p \pmod{pA}\)。将第 1 号的命题 1 应用于 \(\mathbf{J}_k = p^kA\)（\(k \in \mathbf{N}\)）且 \(m = 1\) 的情形，得到同余关系

\[
\Phi_ {n - 1} \left(\sigma \left(a _ {0}\right), \dots , \sigma \left(a _ {n - 1}\right)\right) \equiv \Phi_ {n - 1} \left(a _ {0} ^ {p}, \dots , a _ {n - 1} ^ {p}\right) \pmod{p ^ {n}A},\tag{7}
\]

即

\[
\sigma (u _ {n - 1}) \equiv \Phi_ {n - 1} (a _ {0} ^ {p},..., a _ {n - 1} ^ {p}) \pmod{p ^ {n}A}.\tag{8}
\]

现在根据公式 (2)，关系 \(u_n = \Phi_n(a_0, \ldots, a_n)\) 等价于

\[
u _ {n} = \Phi_ {n - 1} (a _ {0} ^ {p},..., a _ {n - 1} ^ {p}) + p ^ {n} a _ {n}.\tag{9}
\]

引理由此得证。

命题 2. --- 设 A 为一个环。

\begin{enumerate}
\def\labelenumi{\alph{enumi})}
\item
  若 \(p.1_{\mathbf{A}}\) 不是 \(\mathbf{A}\) 中的零因子，则映射 \(\Phi_{\mathbf{A}}\) 是单射。
\item
  若 \(p.1_{\mathbf{A}}\) 在 \(\mathbf{A}\) 中可逆，则映射 \(\Phi_{\mathbf{A}}\) 是双射。
\item
  若 \(\sigma\) 是环 A 的自同态，且满足 \(\sigma(a) \equiv a^p \bmod pA\) 对所有 \(a \in \mathbf{A}\) 成立，则 \(\mathbf{A}'\)，即 \(\Phi_{\mathbf{A}}\) 的像，是 \(\mathbf{A}^{\mathbf{N}}\) 的子环，并且在 \(f_{\mathbf{A}}\) 和 \(v_{\mathbf{A}}\) 下稳定。它是元素 \((u_n)_{n \in \mathbf{N}}\) 的集合，这些元素属于 \(\mathbf{A}^{\mathbf{N}}\) 且满足 \(\sigma(u_n) \equiv u_{n+1} \bmod p^{n+1}A\) 对所有 \(n \in \mathbf{N}\) 成立。
\end{enumerate}

若 \(\boldsymbol{a} = (a_n)_{n \in \mathbb{N}}\) 和 \(\boldsymbol{u} = (u_n)_{n \in \mathbb{N}}\) 是 \(A^{\mathbb{N}}\) 的元素，则根据公式 (2)，关系 \(\Phi_A(\boldsymbol{a}) = \boldsymbol{u}\) 等价于等式

\[
\left\{ \begin{array}{l} u _ {0} = a _ {0}, \\ u _ {n} = \Phi_ {n - 1} \left(a _ {0} ^ {p}, \dots , a _ {n - 1} ^ {p}\right) + p ^ {n} a _ {n} \quad \text {for all} \quad n \geqslant 1. \end{array} \right.\tag{10}
\]

  设 \(\boldsymbol{u} = (u_{n})_{n \in \mathbb{N}}\) 为 \(A^{\mathbb{N}}\) 的元素。当 \(p1_{A}\) 是 A 中的非零因子（分别地，当 \(p1_{A}\) 在 A 中可逆）时，A 中至多存在一个满足等式 (10) 的序列 \((a_{n})_{n \in \mathbb{N}}\)（分别地，恰有一个 A 中满足等式 (10) 的序列 \((a_{n})_{n \in \mathbb{N}}\)），从而得到 a) 和 b)。

 我们证明 \(c)\)。根据引理 2，\(\mathbf{A}'\)，即 \(\Phi_{\mathbf{A}}\) 的像，是元素 \(\boldsymbol{u} = (u_n)_{n \in \mathbb{N}}\) 的集合，这些元素属于 \(\mathbf{A}^{\mathbf{N}}\) 且满足 \(\sigma(u_n) \equiv u_{n+1} \bmod p^{n+1}A\) 对所有 \(n \in \mathbb{N}\) 成立。立即可见 \(\mathbf{A}'\) 是 \(\mathbf{A}^{\mathbf{N}}\) 的子环，在 \(f_{\mathbf{A}}\) 和 \(v_{\mathbf{A}}\) 下稳定。

注记。—设 \(\boldsymbol{a} = (a_n)_{n \in \mathbb{N}}\) 和 \(\boldsymbol{u} = (u_n)_{n \in \mathbb{N}}\) 为 \(A^{\mathbf{N}}\) 中满足 \(\boldsymbol{u} = \Phi_{\mathbf{A}}(\boldsymbol{a})\) 的元素，且 \(m\) 为整数 \(\geqslant 0\)。由 (10) 可推出下列断言：

若 \(u_n\) 在 \(0 \leqslant n \leqslant m\) 时属于 A 的子环 B，且对于每个 \(x \in A\)，关系 \(px \in B\) 蕴含 \(x \in B\)，则 \(a_n\) 在 \(0 \leqslant n \leqslant m\) 时属于 B。

若 A 带有 \(\mathbf{N}\)-分次，\(p1_{A}\) 是 A 中的非零因子，\(d \in \mathbf{N}\)，且 \(u_{n}\) 是次数为 \(dp^{n}\) 的齐次元素，对 \(0 \leqslant n \leqslant m\) 成立，则 \(a_{n}\) 是次数为 \(dp^{n}\) 的齐次元素，对 \(0 \leqslant n \leqslant m\) 成立。

\subsection{多项式的构造}

令 A 为环 Z{[}X, Y{]}，它是具有整数系数、关于两族不定元 \(\mathbf{X} = (\mathbf{X}_{n})_{n \in \mathbb{N}}\) 和 \(\mathbf{Y} = (\mathbf{Y}_{n})_{n \in \mathbb{N}}\) 的多项式环。令 \(\theta\) 为 A 的自同态，由 \(\theta(\mathbf{X}_{n}) = \mathbf{X}_{n}^{p}\) 及 \(\theta(\mathbf{Y}_{n}) = \mathbf{Y}_{n}^{p}\)（对每个 \(n \in \mathbf{N}\)）定义。则 p 不是 A 中的零因子，且 A 中满足 \(\theta(a) \equiv a^{p} \mod p\) 的 a 的集合。A 是包含 X\_n 和 Y\_n 的 A 的子环，因而等于 A。

 根据第 2 号的命题 2 的 a) 和 c)，存在元素 \(\mathbf{S} = (\mathrm{S}_{n})_{n\in\mathbb{N}}\)、\(\mathbf{P} = (\mathrm{P}_{n})_{n\in\mathbb{N}}\)、\(\mathbf{I} = (\mathrm{I}_{n})_{n\in\mathbb{N}}\) 和 \(\mathbf{F} = (\mathrm{F}_{n})_{n\in\mathbb{N}}\)，它们属于 \(\mathrm{A}^{\mathbf{N}}\)，分别由下列等式刻画

\[
\left\{ \begin{array}{l} \Phi_ {\mathbf {A}} (\mathbf {S}) = \Phi_ {\mathbf {A}} (\mathbf {X}) + \Phi_ {\mathbf {A}} (\mathbf {Y}) \\ \Phi_ {\mathbf {A}} (\mathbf {P}) = \Phi_ {\mathbf {A}} (\mathbf {X}) \Phi_ {\mathbf {A}} (\mathbf {Y}) \\ \Phi_ {\mathbf {A}} (\mathbf {I}) = - \Phi_ {\mathbf {A}} (\mathbf {X}) \\ \Phi_ {\mathbf {A}} (\mathbf {F}) = f _ {\mathbf {A}} \big (\Phi_ {\mathbf {A}} (\mathbf {X}) \big). \end{array} \right.\tag{11}
\]

因此，A 中的元素 \(S_{n}\)、\(P_{n}\)、\(I_{n}\) 和 \(F_{n}\) 分别由下列公式刻画（其中 \(n\) 遍历 \(\mathbf{N}\)）：

\begin{enumerate}
\def\labelenumi{(\arabic{enumi})}
\setcounter{enumi}{11}
\tightlist
\item
\end{enumerate}

\[
\Phi_ {n} \left(\mathrm{S} _ {0}, \dots , \mathrm{S} _ {n}\right) = \Phi_ {n} \left(\mathrm{X} _ {0}, \dots , \mathrm{X} _ {n}\right) + \Phi_ {n} \left(\mathrm{Y} _ {0}, \dots , \mathrm{Y} _ {n}\right),\tag{13}
\]

\[
\Phi_ {n} \left(\mathrm{P} _ {0}, \dots , \mathrm{P} _ {n}\right) = \Phi_ {n} \left(\mathrm{X} _ {0}, \dots , \mathrm{X} _ {n}\right) \Phi_ {n} \left(\mathrm{Y} _ {0}, \dots , \mathrm{Y} _ {n}\right),\tag{14}
\]

\[
\Phi_ {n} (\mathrm{I} _ {0}, \dots , \mathrm{I} _ {n}) = - \Phi_ {n} (\mathrm{X} _ {0}, \dots , \mathrm{X} _ {n}),\tag{15}
\]

\[
\Phi_ {n} \left(\mathrm{F} _ {0}, \dots , \mathrm{F} _ {n}\right) = \Phi_ {n + 1} \left(\mathrm{X} _ {0}, \dots , \mathrm{X} _ {n + 1}\right).
\]

给 \(X_{n}\) 和 \(Y_{n}\) 赋予权 \(p^{n}\)，对每个 \(n \in \mathbf{N}\)。由第 2 号的注记推出下列断言：

\begin{enumerate}
\def\labelenumi{\alph{enumi})}
\item
  有 \(\mathbf{S}_n\in \mathbf{Z}[\mathbf{X}_0,\dots ,\mathbf{X}_n,\mathbf{Y}_0,\dots ,\mathbf{Y}_n]\)，且 \(\mathbf{S}_n\) 是权为 \(p^n\) 的等重多项式。
\item
  有 \(\mathbf{P}_n\in \mathbf{Z}[\mathrm{X}_0,\dots ,\mathrm{X}_n,\mathrm{Y}_0,\dots ,\mathrm{Y}_n]\)，且 \(\mathbf{P}_n\) 是权为 \(p^n\) 的等重多项式，分别关于族 \((\mathrm{X}_0,\dots ,\mathrm{X}_n)\) 和 \((\mathrm{Y}_0,\dots ,\mathrm{Y}_n)\)。
\item
  有 \(\mathbf{I}_n\in \mathbf{Z}[\mathrm{X}_0,\dots ,\mathrm{X}_n]\)，且 \(\mathbf{I}_n\) 是权为 \(p^n\) 的等重多项式。
\item
  有 \(\mathbf{F}_n\in \mathbf{Z}[\mathbf{X}_0,\dots ,\mathbf{X}_{n + 1}]\)，且 \(\mathbf{F}_n\) 是权为 \(p^{n + 1}\) 的等重多项式。
\end{enumerate}

公式 (2) 实际上允许逐步确定多项式 \(S_{n}\)、\(P_{n}\)、\(I_{n}\) 和 \(F_{n}\)。

例子。-- 1) 有

\[
\mathbf {S _ {0}} = \mathbf {X _ {0}} + \mathbf {Y _ {0}}
\]

\[
\mathrm{S} _ {1} = \mathrm{X} _ {1} + \mathrm{Y} _ {1} - \sum_ {i = 1} ^ {p - 1} \frac {1}{p} \binom {p} {i} \mathrm{X} _ {0} ^ {i} \mathrm{Y} _ {0} ^ {p - i}.
\]

此外，\(S_{n} - X_{n} - Y_{n}\) 属于环 \(Z[X_{0}, ..., X_{n-1}, Y_{0}, ..., Y_{n-1}]\)。2) 有

\[
\begin{array}{l} \mathbf {P} _ {0} = \mathbf {X} _ {0} \mathbf {Y} _ {0} \\ \mathbf {P} _ {1} = p \mathbf {X} _ {1} \mathbf {Y} _ {1} + \mathbf {X} _ {0} ^ {p} \mathbf {Y} _ {1} + \mathbf {X} _ {1} \mathbf {Y} _ {0} ^ {p}. \end{array}
\]

\begin{enumerate}
\def\labelenumi{\arabic{enumi})}
\setcounter{enumi}{2}
\tightlist
\item
  当 \(p \neq 2\) 时，有 \(I_n = -X_n\)。当 \(p = 2\) 时，有
\end{enumerate}

\[
\mathrm{I} _ {0} = - \mathrm{X} _ {0}
\]

\[
\mathrm{I} _ {1} = - \left(\mathrm{X} _ {0} ^ {2} + \mathrm{X} _ {1}\right)
\]

\[
\mathrm{I} _ {2} = - \mathrm{X} _ {0} ^ {4} - \mathrm{X} _ {0} ^ {2} \mathrm{X} _ {1} - \mathrm{X} _ {1} ^ {2} - \mathrm{X} _ {2}.
\]

\begin{enumerate}
\def\labelenumi{\arabic{enumi})}
\setcounter{enumi}{3}
\tightlist
\item
  有
\end{enumerate}

\[
\mathbf {F} _ {0} = \mathbf {X} _ {0} ^ {p} + p \mathbf {X} _ {1}
\]

\[
\mathrm{F} _ {1} = \mathrm{X} _ {1} ^ {p} + p \mathrm{X} _ {2} - \sum_ {i = 0} ^ {p - 1} \binom {p} {i} p ^ {p - i - 1} \mathrm{X} _ {0} ^ {p i} \mathrm{X} _ {1} ^ {p - i}.
\]

  由于 \(\Phi_n(F_0, \ldots, F_n) \equiv \Phi_n(X_0^p, \ldots, X_n^p) \pmod{p^{n+1}A}\) 对所有 \(n \in \mathbf{N}\) 成立（公式 (2) 和 (15)），根据命题 1(b) 可得 \(F_n \equiv X_n^p \pmod{pA}\) 对所有 \(n \in \mathbf{N}\) 成立。

 注记。--- 令 J 为满足 \(j \geqslant 1\) 的整数集。对于 J 的每个元素 \(j\)，定义多项式 \(\varphi_j\)，它属于 \(\mathbf{Z}[(\mathbf{X}_j)_{j \in \mathbf{J}}]\)，按下式给出

\[
\varphi_ {j} = \sum_ {d} d X _ {d} ^ {j / d},
\]

其中求和遍历 J 中整除 j 的元素。对于每个整数 \(n \geqslant 0\)，有

\[
\varphi_ {p ^ {n}} = \Phi_ {n} (\mathrm{X} _ {p ^ {0}},..., \mathrm{X} _ {p ^ {n}}).
\]

对于每个环 A 及每个元素 \(m\)（属于 \(\mathbf{J}\)），以 \(\varphi_{m}\) 表示从 \(A^{\mathbf{J}}\) 到 A 的映射，它将 \((a_{j})_{j\in\mathbf{J}}\) 送到 \(\varphi_{m}((a_{j})_{j\in\mathbf{J}})\)；以 \(\varphi_{A}\) 或简称 \(\varphi\) 表示从 \(A^{\mathbf{J}}\) 到自身的映射，它将 \(a = (a_{j})_{j\in\mathbf{J}}\) 送到 \((\varphi_{m}(a))_{m\in\mathbf{J}}\)。

令 \(\mathcal{A} = \mathbf{Z}[(X_j)_{j\in\mathbf{J}},(Y_j)_{j\in\mathbf{J}}]\) 为关于两族不定元 \(\mathbf{X} = (X_j)_{j\in\mathbf{J}}\) 和 \(\mathbf{Y} = (Y_j)_{j\in\mathbf{J}}\) 的整数系数多项式环。可以证明（p. 51, 习题 34），\(\mathcal{A}\) 中存在元素

\[
\boldsymbol {s} = (s _ {j}) _ {j \in \mathbf {J}}, \boldsymbol {p} = (p _ {j}) _ {j \in \mathbf {J}} \quad \text{且} \quad \boldsymbol {i} = (i _ {j}) _ {j \in \mathbf {J}},
\]

它们分别由下列等式刻画：

\[
\begin{array}{l} \varphi_ {\mathcal {A}} (s) = \varphi_ {\mathcal {A}} (\mathbf {X}) + \varphi_ {\mathcal {A}} (\mathbf {Y}) \\ \varphi_ {\mathcal {A}} (p) = \varphi_ {\mathcal {A}} (\mathbf {X}) \varphi_ {\mathcal {A}} (\mathbf {Y}) \\ \varphi_ {\mathcal {A}} (i) = - \varphi_ {\mathcal {A}} (\mathbf {X}). \end{array}
\]

\subsection{W(A) 的 Witt 向量环}

设 A 为一个环。若 \(\boldsymbol{a} = (a_n)_{n \in \mathbb{N}}\) 和 \(\boldsymbol{b} = (b_n)_{n \in \mathbb{N}}\) 是 \(\mathbf{A}^{\mathbf{N}}\) 的元素，则以 \(S_{\mathbf{A}}(\boldsymbol{a}, \boldsymbol{b})\)（分别以 \(P_{\mathbf{A}}(\boldsymbol{a}, \boldsymbol{b})\)，以及 \(I_{\mathbf{A}}(\boldsymbol{a})\)），或简称 \(S(\boldsymbol{a}, \boldsymbol{b})\)（分别简称 \(P(\boldsymbol{a}, \boldsymbol{b})\)，以及 \(I(\boldsymbol{a})\)）表示序列 \((S_n(a_0, ..., a_n; b_0, \ldots, b_n))_{n \in \mathbb{N}}\)（分别为 \((P_n(a_0, ..., a_n; b_0, \ldots, b_n))_{n \in \mathbb{N}}\)，以及 \((I_n(a_0, ..., a_n))_{n \in \mathbb{N}}\)）。在公式 (12)、(13) 和 (14) 中，以 \(a_n\) 代替 \(X_n\)、以 \(b_n\) 代替 \(Y_n\)，对所有 \(n \in \mathbb{N}\)，得到等式

\begin{enumerate}
\def\labelenumi{(\arabic{enumi})}
\setcounter{enumi}{15}
\tightlist
\item
\end{enumerate}

\[
\Phi_ {\mathrm{A}} (\mathrm{S} _ {\mathrm{A}} (a, b)) = \Phi_ {\mathrm{A}} (a) + \Phi_ {\mathrm{A}} (b)\tag{17}
\]

\[
\Phi_ {\mathrm{A}} (\mathrm{P} _ {\mathrm{A}} (a, b)) = \Phi_ {\mathrm{A}} (a) \Phi_ {\mathrm{A}} (b)\tag{18}
\]

\[
\Phi_ {\mathbf {A}} \left(\mathrm{I} _ {\mathbf {A}} (\boldsymbol {a})\right) = - \Phi_ {\mathbf {A}} (\boldsymbol {a}).
\]

以 W(A) 表示集合 A \(^{N}\)，并赋予复合律 S \(_{A}\) 和 P \(_{A}\)。

  设 \(\rho : \mathbf{B} \to \mathbf{A}\) 为环同态。以 \(\rho^{\mathbf{N}}\)，或仍以 \(\mathbf{W}(\rho)\)，表示从 \(\mathbf{B}^{\mathbf{N}}\) 到 \(\mathbf{A}^{\mathbf{N}}\) 的映射，它将元素 \(b = (b_n)_{n \in \mathbb{N}}\)（属于 \(\mathbf{B}^{\mathbf{N}}\)）送到 \((\rho(b_n))_{n \in \mathbb{N}}\)。由定义立即得到

\begin{enumerate}
\def\labelenumi{(\arabic{enumi})}
\setcounter{enumi}{18}
\tightlist
\item
\end{enumerate}

\[
\mathrm{W} (\rho) \circ \mathrm{S} _ {\mathrm{B}} = \mathrm{S} _ {\mathrm{A}} \circ (\mathrm{W} (\rho) \times \mathrm{W} (\rho))\tag{20}
\]

\[
\mathrm{W} (\rho) \circ \mathrm{P} _ {\mathrm{B}} = \mathrm{P} _ {\mathrm{A}} \circ (\mathrm{W} (\rho) \times \mathrm{W} (\rho))\tag{21}
\]

\[
\mathrm{W} (\rho) \circ \mathrm{I} _ {\mathrm{B}} = \mathrm{I} _ {\mathrm{A}} \circ \mathrm{W} (\rho)\tag{22}
\]

\[
\rho^ {\mathbf {N}} \circ \Phi_ {\mathrm{B}} = \Phi_ {\mathrm{A}} \circ \mathrm{W} (\rho).
\]

 引理 3. --- 设 A 为一个环。存在满射环同态 \(\rho: \mathbf{B} \to \mathbf{A}\)，其中 \(\mathbf{B}\) 为满足下列条件的环：p 不是 \(\mathbf{B}\) 中的零因子，且存在自同态 \(\sigma\)，它是 \(\mathbf{B}\) 的自同态，使得 \(\sigma(b) \equiv b^p\) mod. \(p\) . \(\mathbf{B}\) 对每个 \(b \in \mathbf{B}\) 成立。

  事实上，只需置 \(\mathbf{B} = \mathbf{Z}[(X_{a})_{a \in A}]\)，取 \(\sigma\) 为 \(\mathbf{B}\) 的自同态，由 \(\sigma(X_{a}) = X_{a}^{p}\) 对每个 \(a \in A\) 定义，并取 \(\rho\) 为从 \(\mathbf{B}\) 到 A 的同态，由 \(\rho(X_{a}) = a\) 对每个 \(a \in A\) 定义即可。

定理 1. --- a) 设 A 为一个（交换）环。赋予 \(S_A\) 加法和 \(P_A\) 乘法后，\(W(A)\) 是一个（交换）环。加法单位元是序列 \(0_A\)，其所有项均为零；乘法单位元是序列 \(1_A\)，除指标为 0 且等于 \(1_A\) 的项外，其余各项均为零。元素 \(a\)（属于 \(W(A)\)）的加法逆元是 \(I_A(a)\)。

\begin{enumerate}
\def\labelenumi{\alph{enumi})}
\setcounter{enumi}{1}
\item
  设 \(\rho: \mathbf{B} \to \mathbf{A}\) 为环同态。则 \(\mathbf{W}(\rho): \mathbf{W}(\mathbf{B}) \to \mathbf{W}(\mathbf{A})\) 是环同态。
\item
  设 A 为一个环。映射 \(\Phi_{\mathbf{A}}\) 是从 \(W(A)\) 到乘积环 \(\mathbf{A}^{\mathbf{N}}\) 的环同态。特别地，对于每个 \(n \in \mathbf{N}\)，映射 \(\Phi_{n}: a \mapsto \Phi_{n}(a_{0}, \ldots, a_{n})\) 是从 W(A) 到 A 的环同态。
\end{enumerate}

  由公式 (16)、(17)、(19) 和 (20)，只需证明断言 \(a\)。设 \(\rho: \mathbf{B} \to \mathbf{A}\) 为满足引理 3 条件的环同态。令 \(\mathbf{B}'\) 为 \(\mathbf{B}^{\mathbf{N}}\) 的子环，由元素 \((b_n)_{n \in \mathbf{N}}\) 构成，且满足 \(\sigma(b_n) \equiv b_{n+1} \pmod{p^{n+1}\mathbf{B}}\) 对所有 \(n \in \mathbf{N}\) 成立。根据第 2 号的命题 2，\(\Phi_{\mathbf{B}}\) 诱导出 \(\Phi_{\mathbf{B}}'\)，它是从 W(B) 到 \(\mathbf{B}'\) 的双射。结合公式 (16) 至 (18) 以及关系 \(\Phi_n(\mathbf{0}_{\mathbf{B}}) = 0\) 和 \(\Phi_n(\mathbf{1}_{\mathbf{B}}) = \mathbf{1}_{\mathbf{B}}\)（\(n \in \mathbf{N}\)），通过结构迁移可见 W(B) 是一个环，其加法单位元为 \(\mathbf{0}_{\mathbf{B}}\)，乘法单位元为 \(\mathbf{1}_{\mathbf{B}}\)，且 b 的相反元等于 \(\mathbf{I}_{\mathbf{B}}(b)\)。

映射 W(ρ): W(B) → W(A) 是满射。根据公式 (19) 和 (20)，与映射 W(ρ) 相关的 W(B) 上的等价关系 R 与 W(B) 的环结构相容。由于 W(ρ) 诱导出从商环 W(B)/R 到 W(A) 的双射 Ψ，且与加法和乘法相容，断言 a) 由结构迁移得出。

定义 1. --- 设 A 为一个环。称环 \(W(A)\) 为系数在 A 中的 Witt 向量环。

对于元素 \(\boldsymbol{a}\)（属于 \(W(A)\)）及 \(n\)（属于 \(\mathbf{N}\)），元素 \(\Phi_n(\boldsymbol{a}) = \Phi_n(a_0, \ldots, a_n)\) 有时称为指标 \(n\) 的 \(\boldsymbol{a}\) 的幽灵分量。

注记。--- 重新采用第 3 号注记的记号。设 A 为一个环。若 \(a\) 和 \(b\) 是 \(\mathbf{A}^{\mathbf{J}}\) 中的元素，且 \(r = (r_{j})_{j\in\mathbf{J}}\) 是 \(\mathcal{A}^{\mathbf{J}}\) 中的元素，则以 \(r_{\mathrm{A}}(a,b)\) 表示元素 \((r_{j}(a,b))_{j\in\mathbf{J}}\)，它属于 \(\mathbf{A}^{\mathbf{J}}\)。以 \(U(A)\) 表示集合 \(\mathbf{A}^{\mathbf{J}}\)，并赋予其复合律 \(s_{\mathrm{A}}\) 和 \(p_{\mathrm{A}}\)。可以证明（p. 52, 习题 35），赋予加法 \(s_{\mathrm{A}}\) 和乘法 \(p_{\mathrm{A}}\) 后，\(U(A)\) 是一个（交换）环；称其为 A 的万有 Witt 环。加法单位元是 \(U(A)\) 中所有分量均为零的元素；乘法单位元是 \(U(A)\) 中除指标为 1 且等于 \(1_{\mathrm{A}}\) 的分量外其余分量均为零的元素；元素 \(a\)（属于 \(U(A)\)）的相反元是 \(i_{\mathrm{A}}(a)\)。映射 \(\varphi_{\mathrm{A}}\) 是从 \(U(A)\) 到乘积环 \(\mathbf{A}^{\mathbf{J}}\) 的环同态。

设 \(\rho: \mathbf{B} \to \mathbf{A}\) 为环同态；以 \(\mathrm{U}(\rho)\) 表示从 \(\mathbf{B}^{\mathbf{J}}\) 到 \(\mathbf{A}^{\mathbf{J}}\) 的映射，它将元素 \((b_j)_{j \in \mathbf{J}}\)（属于 \(\mathbf{B}^{\mathbf{J}}\)）送到元素 \((\rho(b_j))_{j \in \mathbf{J}}\)（属于 \(\mathbf{A}^{\mathbf{J}}\)）。可以证明（同上）\(\mathrm{U}(\rho)\) 是从 \(\mathrm{U}(\mathrm{B})\) 到 \(\mathrm{U}(\mathrm{A})\) 的环同态。
